%%%%%%%%%%%%%%%%%%%%%%%%%%%%%%%%%%%%%%%%%%%%%%%%%%%%%%%%%%%%%%%%%%%%%%%%%%%%%%%%
% NDSS 2026 submission --- IEEEtran conference format.
%
% "From PKI Standards to Compliance-Checking Code: An End-to-End Framework
%  for Rule Extraction, Lintability Determination, and Verifiable Code
%  Generation."
%
% Packages and the author block live in the project's own packages.tex and
% authors.tex (\input below); this file only adds the paper-specific macros,
% theorem environments, and document structure. The section content
% (0_abstract, 1_introduction ... 10_conclusion, appendix, algorithms/) is the
% same as the USENIX version, included unchanged --- no typewriter markup was
% added and existing bold was left as-is.
%
% Build with the IEEEtran.cls used by the NDSS template:
%     pdflatex main ; bibtex main ; pdflatex main ; pdflatex main
%%%%%%%%%%%%%%%%%%%%%%%%%%%%%%%%%%%%%%%%%%%%%%%%%%%%%%%%%%%%%%%%%%%%%%%%%%%%%%%%

\documentclass[conference]{IEEEtran}

% Packages + comment/shorthand macros (project file, merged: algorithmic ->
% algpseudocode, and amsthm/dsfont/booktabs/multirow/stfloats added there).
\usepackage{cite}
\usepackage{amsmath,amssymb,amsfonts}
\usepackage{algpseudocode} % algorithmicx syntax inside custom algorithm listing floats
\usepackage{array}
\usepackage{graphicx}
\usepackage{textcomp}
\usepackage{xcolor}
\usepackage{xurl}

% Additional packages required by this paper's content:
\usepackage{amsthm}      % theorem/proposition environments (appendix)
\usepackage{dsfont}      % \mathds{1} indicator function (\ind)
\usepackage{booktabs}    % \toprule / \midrule / \bottomrule in tables
\usepackage{tabularx}    % width-constrained tables (X column absorbs slack)
\usepackage{multirow}    % multirow table cells
\usepackage{stfloats}    % full-width (starred) floats at page bottom (modern; avoids deprecated fixltx2e)
\usepackage{enumitem}    % vertical spacing of itemize/enumerate (spacing only; margins untouched)

% IEEEtran applies \scshape to table-caption text, but the Times companion
% font used here has no such shape. Keep IEEEtran's caption layout and make
% that one local request upright instead.
\makeatletter
\let\NDSSoriginalmakecaption\@makecaption
\long\def\@makecaption#1#2{%
  \begingroup
  \let\scshape\upshape
  \NDSSoriginalmakecaption{#1}{#2}%
  \endgroup
}
\makeatother

% Character protrusion only. expansion=false is deliberate: font expansion would
% scale glyph widths, which the NDSS/IEEEtran template forbids ("do not use
% packages that alter fonts"). Protrusion only shifts punctuation optically into
% the margin edge and leaves every glyph, the typeface, and the body font size
% unchanged.
\usepackage[protrusion=true,expansion=false]{microtype}

% Tighten only the vertical spacing of lists. leftmargin is deliberately left at
% the class default so no length controlling horizontal margins is touched.
\setlist[itemize]{topsep=2pt,partopsep=0pt,parsep=0pt,itemsep=2pt}
\setlist[enumerate]{topsep=2pt,partopsep=0pt,parsep=0pt,itemsep=2pt}

\newif\ifshowcomments
\showcommentsfalse
\newcommand{\zmm}[1]{\ifshowcomments{\color{purple}{\{zmm: #1\}}}\fi}
\newcommand{\zmmmod}[1]{{\color{teal}{#1}}}
\newcommand{\revise}[1]{{\color{black}{#1}}}
\newcommand{\todo}[1]{{\color{red}{\{Todo: #1\}}}}

\newcommand{\ignore}[1]{}
\newcommand{\etal}{et al. }
\newcommand{\ie}{i.e., }
\newcommand{\eg}{e.g., }

% Algorithm listings use algorithmicx inside ordinary IEEE floats, but keep an
% independent counter so they are referenced as algorithms rather than figures.
\newcounter{algorithm}
\renewcommand{\thealgorithm}{\arabic{algorithm}}


% IEEE: restore multiline-equation page breaks (amsmath sets this to 10000).
\interdisplaylinepenalty=2500

% Extra float slots --- this paper has many full-width table*/figure* floats.
\extrafloats{64}

\pagestyle{plain}   % show page numbers after the title page (NDSS submission)

%-------------------------------------------------------------------------------
% Theorem-like environments (amsthm is loaded in packages.tex).
%-------------------------------------------------------------------------------
\newtheorem{theorem}{Theorem}
\newtheorem{proposition}{Proposition}
\newtheorem{definition}{Definition}

%-------------------------------------------------------------------------------
% Paper-specific math shorthands used across the section files.
%-------------------------------------------------------------------------------
\newcommand{\ind}{\mathds{1}}                 % indicator function
\newcommand{\code}{\mathrm{Code}}
\newcommand{\ir}{\mathrm{IR}}
\newcommand{\spec}{\mathrm{Spec}}
\newcommand{\CodeEqIR}{\ensuremath{\code\equiv\ir}}
\newcommand{\CodeEqSpec}{\ensuremath{\code\equiv\spec}}
\newcommand{\Rkw}{\mathcal{R}_{\mathrm{kw}}}
\newcommand{\smech}{\sigma_{\mathrm{mech}}}
\newcommand{\Tv}{\mathcal{T}_{\mathcal{V}}}
\newcommand{\Lrec}{\mathcal{L}_{\mathrm{recall}}}
\newcommand{\Lcode}{\mathcal{L}_{\mathrm{code}}}
\newcommand{\Ltot}{\mathcal{L}_{\mathrm{total}}}
\newcommand{\Nviol}{N_{\mathrm{viol}}}
\newcommand{\RL}{\mathcal{R}_L}
\newcommand{\RLcov}{\mathcal{R}_L^{\mathrm{cov}}}
\newcommand{\RLunc}{\mathcal{R}_L^{\mathrm{uncov}}}
\newcommand{\RN}{\mathcal{R}_N}
\newcommand{\RU}{\mathcal{R}_U}

%-------------------------------------------------------------------------------
% Fixed-width table column types (L/C/R) used in the methodology/appendix tables
% (array is loaded in packages.tex). \ie/\eg and the review macros come from
% packages.tex, so they are not redefined here.
%-------------------------------------------------------------------------------
\newcolumntype{L}[1]{>{\raggedright\arraybackslash}p{#1}}
\newcolumntype{C}[1]{>{\centering\arraybackslash}p{#1}}
\newcolumntype{R}[1]{>{\raggedleft\arraybackslash}p{#1}}

%-------------------------------------------------------------------------------
% \cicasfig{path}{fallback caption text}
% Uses the real figure if figures/<name> exists; otherwise draws a labeled
% placeholder box so the project always compiles (even before the PNGs are added).
%-------------------------------------------------------------------------------
\newcommand{\cicasfig}[2]{%
  \IfFileExists{#1}{%
    \includegraphics[width=\linewidth]{#1}%
  }{%
    \fbox{\parbox[c][3.2cm][c]{0.94\linewidth}{\centering
      \textit{[figure placeholder --- add:}\\[0.2em]
      {\ttfamily\footnotesize \detokenize{#1}}\textit{]}\\[0.5em]
      \footnotesize #2}}%
  }%
}

%-------------------------------------------------------------------------------
\begin{document}
%-------------------------------------------------------------------------------

\title{Compiling PKI Standards into Compliance-Checking Code}

% Author block (project file --- fill in the real names/affiliations there).
% Intentionally empty for double-blind review.
\author{}


% NDSS 2026 pubid banner (from bare_conf_NDSS2026.tex).
\IEEEoverridecommandlockouts
\makeatletter\def\@IEEEpubidpullup{6.5\baselineskip}\makeatother
\IEEEpubid{\parbox{\columnwidth}{
    Network and Distributed System Security (NDSS) Symposium 2026\\
    23 - 27 February 2026 , San Diego, CA, USA\\
    ISBN 979-8-9919276-8-0\\
    https://dx.doi.org/10.14722/ndss.2026.[23$|$24]xxxx\\
    www.ndss-symposium.org
}
\hspace{\columnsep}\makebox[\columnwidth]{}}

\maketitle

%-------------------------------------------------------------------------------
\begin{abstract}
%-------------------------------------------------------------------------------
Web PKI trust depends on correctly issued certificates. Static certificate
linters such as zlint enforce this trust, but PKI experts must still implement
and maintain compliance checks as normative sources evolve. This manual process
is increasingly difficult to scale as pre-issuance linting becomes mandatory and
certificate lifetimes shorten. The ecosystem lacks a systematic way to translate
natural-language normative text into lint code, including a principled method to
decide which rules are observable from one certificate alone. We present an
end-to-end framework that extracts normative rules from Web PKI specifications,
classifies their lintability, and generates auditable lint code. A dual-layer
extraction pipeline combines deterministic candidate recall with
schema-constrained large-language-model parsing into an intermediate
representation. A four-condition decision rule then classifies lintability from
the representation rather than from keyword heuristics. Restricted code
generation maps supported rules to a closed, typed DSL tree that a deterministic
renderer converts to Go lint code, with a build-gated LLM fallback for unresolved
cases. On RFC~5280 and the CA/Browser Forum Baseline Requirements, the framework
recalls 2,080 candidate rules, of which 1,567 are genuine normative rules. Only
260 (16.6\%) are lintable as single-certificate static checks; existing zlint
lints fully cover 165. For the remaining 95 rules, the generator produces 93
compilable lints, 79 of which pass a unanimous code--specification synonymy gate.
When added to zlint, these lints confirm five certificate-problem types across
the evaluated corpora, including findings missed by native zlint in a scan of
111,118 deployed certificates.

\end{abstract}

\IEEEpeerreviewmaketitle

%-------------------------------------------------------------------------------
\section{Introduction}
\label{sec:intro}
%-------------------------------------------------------------------------------

The Public Key Infrastructure (PKI) forms the trust foundation of the Internet, enabling identity authentication and encrypted communication through digital certificates issued by Certificate Authorities (CAs).
Correct certificate issuance is therefore critical to Web PKI integrity, because certificates that violate technical requirements directly undermine browser trust and weaken the reliability of secure communication.

Recent incidents demonstrate that these failures are not merely theoretical.
In 2024, both Chrome and Mozilla announced that they would distrust Entrust as a publicly trusted CA after a series of unresolved compliance incidents~\cite{mozilla2024entrust,chrome2024entrust}.
Earlier, in 2020, Let's Encrypt revoked approximately three million certificates after a CAA-validation defect~\cite{letsencrypt2020caa}.
% These incidents have already translated into ecosystem-wide policy changes.
\zmmmod{In parallel, the Web PKI ecosystem has moved toward stronger automated compliance controls.}
The CA/Browser Forum mandated pre-issuance linting for all publicly trusted CAs beginning March~15,~2025~\cite{cabf-sc75-2024}.
Subsequently, Ballot~SC-081v3~\cite{cabf-sc081v3} introduced a progressive reduction of TLS certificate validity periods to 47~days, substantially increasing certificate issuance frequency and, \zmmmod{consequently, the operational importance of scalable automated compliance checking.}

Currently, certificate compliance checking relies primarily on linting tools such as zlint~\cite{kumar2018tracking}, pkilint~\cite{pkilint}, certlint~\cite{certlint}, and x509lint~\cite{x509lint}, which are widely integrated into Certificate Transparency (CT)~\cite{stark2019does} monitoring and CA issuance workflows.
% However, these tools depend heavily on manual certificate-lint engineering, since each certificate lint must be individually implemented and continuously maintained by PKI experts as Web PKI normative sources are revised.
% As certificate issuance volume and the complexity of Web PKI normative sources continue to grow, this approach becomes increasingly difficult to scale.
\zmmmod{However, translating normative requirements into executable checking lints remains largely manual.
As Web PKI standards evolve, experts must identify, interpret, and encode applicable requirements as checks, making the process increasingly difficult to scale with the growth of certificate issuance and the normative ecosystem.}


% More fundamentally, the PKI ecosystem lacks a systematic framework for translating normative text into compliance logic enforceable via certificate lints.
% This gap matters as linting becomes increasingly critical to PKI governance, because Web PKI normative sources still do not distinguish normative rules that admit certificate lints from those that require evidence not contained in the certificate.
\zmmmod{
More fundamentally, existing linting tools begin only after a requirement has already been understood and deemed suitable for certificate-level checking. 
What is missing is a systematic mechanism that starts from the normative sources themselves and determines which requirements can be translated into executable certificate lints.
This distinction is essential because \textit{normative does not imply lintable}: many mandatory requirements depend on operational procedures, external state, or runtime behavior that cannot be determined from certificate contents alone.
Consequently, the Web PKI ecosystem lacks a systematic framework from normative text to enforceable certificate-lint logic.
}

\noindent \textbf{Research goal and challenges.}
We aim to develop a general framework that automatically extracts normative rules from Web PKI normative sources, determines whether each normative rule can be enforced through certificate linting, and generates the corresponding lint code.
This task presents three key challenges.

% First, normative-rule extraction must be accurate. 
\zmmmod{First, the semantics of Web PKI normative rules are often non-local.}
Normative sources contain numerous cross-section and cross-document references, inheritance relations, and strengthened constraints.
For example, CABF~BR~\S7.1 inherits normative requirements from RFC~5280, while CABF~BR~\S7.1.2.7.12 strengthens RFC~5280~\S4.2.1.6 by requiring the subjectAltName extension in every Subscriber Certificate.
\ignore{so isolated analysis of a single source is insufficient. 
Large Language Models (LLMs) offer a promising basis for this semantic understanding. To retain an inspectable link between the source material and later analysis, our design uses source-grounded retrieval and a structured intermediate representation rather than an end-to-end black-box judgment.}
\zmmmod{Correctly interpreting such a rule therefore requires recovering its relevant normative context rather than analyzing an isolated sentence or section.}

% Second, lintability determination is itself non-trivial.\footnote{Lintability in this study denotes the extent to which a normative rule can be implemented as a certificate lint without dependence on external context or runtime behavior.}
\zmmmod{Second, determining whether a normative rule is lintable is a non-trivial semantic problem.\footnote{In this study, a normative rule is considered lintable when its compliance can be determined using the information available to a certificate lint, without requiring external operational evidence or runtime behavior.}}
The presence of RFC~2119 keywords such as MUST, SHALL, or SHOULD does not imply that a normative rule can be translated into a deterministic certificate lint.
Many keyword-bearing rules depend on external context or runtime behavior. 
We therefore need an explicit, checkable characterization of lintability rather than a keyword heuristic.

% Third, code generation must be trustworthy~\cite{liu2023evalplus}. 
% An LLM-to-code workflow without an explicit output-space restriction does not enforce the field selections, constants, or program structures represented by a normative rule. Our design instead confines code generation to a closed DSL-tree space and applies structural and build checks before a candidate enters downstream verification.
\zmmmod{Third, generated lint code must faithfully preserve the semantics of the originating normative rule.
An unconstrained LLM-to-code workflow may produce plausible or compilable programs while selecting incorrect certificate fields, constants, predicates, or control structures~\cite{liu2023evalplus}.
Thus, automated lint generation requires a restricted and auditable output space in which syntactic validity and semantic correspondence can be systematically checked. 
}

\noindent \textbf{Our study.}
% We design a pipeline that automatically extracts lintable normative rules from Web PKI normative sources, determines their lintability, and generates auditable lint code, addressing the three challenges in turn.
\zmmmod{We design an end-to-end framework for compiling Web PKI normative text to compliance-checking lints.
It automatically extracts normative rules from Web PKI normative sources, determines their lintability, and generates auditable lint code, addressing the three challenges.
A central design principle is to use LLMs only for bounded semantic interpretation while making their inputs, outputs, and downstream effects structurally constrained and auditable.
}

To address the first challenge, we develop the Dual Layer Extraction Pipeline (DLEP). 
We construct a knowledge graph over Web PKI normative sources that captures sections, definitions, certificate fields, and cross-document relations, 
% and we design a deterministic retrieval mechanism that assembles context via bounded traversal over explicit reference edges, giving traceable context for extraction. 
\zmmmod{and use deterministic traversal to retrieve the relevant source context.}
% DLEP then extracts normative rules in two layers. 
% Layer~1 performs deterministic normative-rule recall using keyword filtering and regex matching to identify candidate rules.
% Layer~2 employs an LLM for constrained semantic parsing of these candidates into a JSON-serialized Intermediate Representation (IR).
% This confines the LLM to schema-bounded interpretation of source-grounded context. 
% The resulting IR includes fields for source spans, evidence text, and retrieved context nodes.
\zmmmod{DLEP then extracts normative rules in two layers: 
the first uses deterministic keyword and regex matching to recall candidate rules, 
while the second employs an LLM to parse them into a schema-constrained Intermediate Representation (IR) that preserves links to the supporting source text and retrieved context.
This design grounds semantic interpretation in structured source evidence, avoiding reliance on end-to-end black-box judgment.
}

% To address the second challenge, we operationalize lintability through a four-condition framework in \S\ref{sec:method-lintability}, evaluated over IR attributes, turning the lintability determination, given a fixed IR, into a traceable deterministic computation rather than a keyword heuristic.
\zmmmod{To address the second challenge, we determine lintability through a four-condition framework that evaluates whether the compliance decision can be expressed using information available to a certificate lint.
Given a fixed IR, the framework turns lintability determination into a traceable and reproducible computation rather than a heuristic based on keywords.}

% To address the third challenge, we restrict code generation to a finite, closed DSL-tree space. Each normative rule's checking logic is a tree in this space that a deterministic function renders into Go lint code. The generator first attempts deterministic synthesis; when that path yields no tree, an LLM fallback may propose a tree from the same closed space. Every candidate is gated by a requirement that it render and pass go build, so out-of-space and non-compiling candidates are rejected before downstream semantic evaluation.
\zmmmod{To address the third challenge, we restrict generated checking logic to a closed DSL-tree representation that admits deterministic rendering into Go lint code.
Both deterministic synthesis and LLM-assisted generation operate within this restricted representation, and generated candidates are subjected to structural, build, and downstream semantic checks before acceptance.
The goal is not merely to generate compilable code, but to constrain generation so that the relationship between a normative rule and its executable check remains inspectable.}

\noindent \textbf{Evaluation and Analysis.}
We evaluate the full framework on two central Web PKI normative sources RFC~5280 and CABF~BR.
Of 2,080 keyword-recalled candidate records, 1,567 are retained as normative rules by the pipeline's adjudication procedure, of which 260 (16.6\%) are determined to be lintable.
Existing zlint lints cover 165 of these rules, leaving 95 lintable normative rules (36.5\%) without full zlint coverage. 
% and revealing a substantial gap between deployed tooling and the normative sources.
% For the uncovered rules, we report the code-generation rate using the
% GENERIC-participating count, which is 76/95 (80.0\%), consisting of 72/95 (75.8\%)
% GENERIC-only generated lints and 4/95 (4.2\%) lints that mix GENERIC and
% NON\_GENERIC atoms.
% Sixty-four of these pass the final unanimous code--normative-rule semantic-equivalence gate, a rate of 64/76 (84.2\%).
\zmmmod{For these uncovered rules, our framework generates candidate lints for 76 of 95 rules (80.0\%), and 64 of the 76 generated candidates (84.2\%) pass the final semantic-equivalence validation.}

We further perform two independent external validations, targeting the two major halves of the framework.
Against the zlint maintainers' manually annotated CABF~BR mapping sheet~\todo{add citation}, our normative-rule extraction and lintability determination show substantial agreement with expert annotations across both available versions.
We additionally integrate the accepted generated lints into zlint and apply them to the evaluated certificate corpora.
This process surfaces five confirmed classes of findings, including violations in deployed certificates that are not detected by native zlint, demonstrating that the uncovered normative rules correspond to practical gaps in existing tooling.

\noindent\textbf{Contributions.}
This paper makes three contributions.
\begin{itemize}
    \item 
    % We design DLEP, a dual-layer pipeline that automatically extracts normative rules from natural-language Web PKI normative sources, confines the LLM to schema-bounded semantic parsing, and produces source-linked IR records for inspection.
    We present DLEP, a source-grounded framework for translating natural-language Web PKI normative text into structured, auditable rule representations, combining deterministic context retrieval with schema-constrained semantic parsing.
    \item We propose a four-condition operational criterion for determining whether an extracted normative rule can be enforced as a certificate lint.
    \item 
    % We develop a restricted code-generation method that confines each normative rule's checking logic to a finite, closed DSL-tree space, yielding deterministic lint-code rendering with explicit audit evidence and a build-gated LLM fallback.
    We develop a constrained normative-rule-to-lint synthesis approach that restricts checking logic to a closed DSL-tree representation, enabling deterministic code rendering, explicit audit evidence, and systematic structural and semantic validation.
\end{itemize}

%-------------------------------------------------------------------------------
\section{Background and Related Work}
\label{sec:background}
%-------------------------------------------------------------------------------

This section reviews the domain background and the prior work that motivate our design. We describe the Web PKI normative-source stack and its multi-tier constraints on certificate format, summarize the certificate-linting tools whose limitations motivate this study, and survey related work on specification extraction, knowledge-graph retrieval, and LLM output control.

\subsection{X.509 Certificate and Web PKI Normative Sources}
\label{sec:bg-pki}

X.509 certificates are the primary data representation in the Web PKI. Their structure and validation semantics are defined across layered normative sources and policy frameworks. IETF~RFC~5280 defines the foundational layer, including the certificate ASN.1 structure, field-level semantics, and the path-validation algorithm. Above it, the CA/Browser Forum Baseline Requirements (BR) impose operational and configuration constraints on commercial CAs that issue publicly trusted TLS certificates. The ETSI~EN\,319 series adds EU-specific configuration profiles aligned with the eIDAS trust framework, and the Mozilla Root Store Policy (MRSP) occupies the browser-policy layer with its own technical requirements. Together, these four normative sources form the Web PKI normative-source stack. The layers cross-reference one another and sometimes impose local overrides. Chapter~7 of CABF~BR, for example, references RFC~5280 directly while tightening individual field-level constraints.

A certificate is a signed ASN.1 object whose core payload is the tbsCertificate structure. This payload contains identity and validity metadata, including version, serialNumber, signature, issuer, validity, subject, and subjectPublicKeyInfo. Version~3 certificates further carry an extension sequence, which is where many security-relevant Web PKI constraints reside. Common examples include subjectAltName for DNS names and email names, basicConstraints for CA status and path-length limits, keyUsage and extendedKeyUsage for allowed cryptographic purposes, certificatePolicies for policy identifiers, and authority-information or revocation pointers such as AIA and CRL distribution points. These fields are encoded as DER bytes, but normative rules are usually written as natural-language statements over semantic field names, profiles, and certificate roles.

\subsection{Certificate Linting}
\label{sec:bg-linting}

Web PKI certificate issuance compliance is checked primarily with certificate-linting tools. These certificate-linting tools inspect the encoded contents of individual certificates and do not require external validation state. zlint~\cite{kumar2018tracking} is the most widely deployed certificate linter. pkilint~\cite{pkilint}, certlint~\cite{certlint}, and x509lint~\cite{x509lint} extend this ecosystem further.

zlint organizes checks as individually registered lints with metadata such as citation, effective date, and severity; each lint implements a predicate over parsed certificate structures and is commonly exercised through pass/error test certificates.

This design has two important implications for our work. First, a certificate lint is not a direct copy of a normative source sentence. It is an engineering artifact that embeds a chosen field path, a predicate, a severity level, and often a scope condition. Second, the absence of a certificate lint does not necessarily mean that no normative rule exists; it may mean that the normative rule has not been recognized as lintable or has not yet been implemented.

\subsection{Related Work}
\label{sec:bg-related}

PolicyLint~\cite{andow2019policylint} and PrivacyFlash~Pro~\cite{zimmeck2021privacyflash}
analyze privacy text~\cite{harkous2018polisis,amaral2022ai}, Hassani et al.~\cite{hassani2024rethinking} use LLMs for
legal-compliance analysis~\cite{sun2025compliance,masoudifard2024leveraging,chung2025graphcompliance}, and N-Check~\cite{feng2024analyzing} formalizes
normative rules as FOL$^*$ formulas for well-formedness analysis. These systems
extract or formalize obligations, but do not map Web PKI issuance rules to
DER-encoded certificate fields or determine their single-certificate lintability.

Specification analysis covers protocol states, formats, and code~\cite{pacheco2022automated,yen2021sage}.
PROSPER~\cite{sharma2023prosper}, ParCleanse~\cite{zheng2025validating}, and
SpecGPT~\cite{zhang2025automated} extract RFC state machines, protocol formats,
and 3GPP state-machine behavior. Wu et al.~\cite{wu2025uncovering} align RFC text
with kernel code using GraphRAG-style context~\cite{lewis2020retrieval,edge2024local,peng2025graph},
while RFCAUDIT~\cite{zheng2025rfcaudit} uses hierarchical code-semantic indexing
and retrieval-guided LLM analysis to find inconsistencies between RFC rules and
protocol implementations. SPEC2CODE~\cite{wang2025spec2code} maps fine-grained
RFC rules to functions, ProtocolGuard~\cite{song2026protocolguard} combines rule
extraction with static and dynamic verification of existing implementations, and
LLMs Unleashed~\cite{long2026llms} generates complete protocol implementations.
Unlike our pipeline, these systems do not jointly provide exhaustive rule recall,
schema-constrained IR parsing, deterministic lintability determination, and
restricted lint-predicate generation.

Output-control and verification work addresses complementary layers.
XGrammar~\cite{dong2025xgrammar} and CRANE~\cite{banerjee2025crane} constrain LLM
outputs through grammars, although constrained decoding can impair reasoning
quality~\cite{schall2025hidden,park2024grammar}. ARMOR~\cite{debnath2024armor} verifies the
RFC~5280 path-validation algorithm in Agda, while EverParse~\cite{ramananandro2019everparse}
and Nail~\cite{bangert2014nail} generate verified parsers and serializers from
declarative format grammars. They improve output form, algorithm assurance~\cite{debnath2021reengineering}, or
format handling rather than translating natural-language issuance rules into
auditable certificate lints.

Certificate studies provide the closest empirical context~\cite{durumeric2013analysis,georgiev2012most,amann2017mission,ma2021tracing,zhang2021rusted}. Kumar et
al.~\cite{kumar2018tracking} measured misissuance at ecosystem scale, and Zhang et
al.~\cite{zhang2025unicert} identified Unicode-related issuance and parsing
non-compliance, but both begin with known lints or defect classes.
Frankencerts~\cite{brubaker2014frankencerts} uses adversarial certificates for
differential testing, and SymCerts~\cite{chau2017symcerts} uses symbolic execution
to expose validator noncompliance. These systems exercise validation code rather
than derive lintable issuance rules and lint predicates from normative text.

%-------------------------------------------------------------------------------
\section{Rule Extraction and Lintability Determination}
\label{sec:methodology}
%-------------------------------------------------------------------------------

This section presents DLEP, our framework for Web PKI normative rule extraction and lintability analysis.
We describe the overall system architecture and the methods used in the key components.

\subsection{Architecture Overview}
\label{sec:method-overview}

Web PKI normative sources serve as the input to the framework. Through keyword recall and constrained extraction, candidate rules are converted into a structured representation. As shown in Fig.~\ref{fig:architecture}, the framework contains four modules.

(1)~Knowledge graph construction. To handle scattered definitions, inherited requirements, and cross-document references, this module organizes the collected sources into a traceable source context for rule extraction.

(2)~Deterministic context retrieval. To prevent the LLM from inventing unsupported context, this module supplies source-grounded context for each target section in a reproducible way.

(3)~DLEP. To balance completeness and precision, this module first recalls candidate normative rules; for precise extraction, the second layer uses an LLM to convert them into a structured representation.

(4)~Four-condition lintability framework. To determine whether a rule can be implemented as a static lint check, we design a four-condition lintability framework.

\begin{figure*}[tb]
  \centering
  \cicasfig{figures/fig01_architecture.png}{Architecture of the Dual-Layer
    Extraction Pipeline (DLEP), showing knowledge-graph construction, deterministic
    context retrieval, dual-layer rule extraction, and the four-condition
    lintability framework.}
  \caption{Architecture of Dual-Layer Extraction Pipeline (DLEP).}
  \label{fig:architecture}
\end{figure*}

\subsection{PKI Knowledge Graph Builder}
\label{sec:method-kg-builder}

The PKI knowledge graph builder is the offline component that converts heterogeneous Web PKI normative sources into the retrieval substrate used by DLEP.
It normalizes documents into a hierarchical structure with stable identifiers for specifications, sections, headings, and text units.
It then extracts explicit references such as section citations, RFC references, and policy cross-links, and links them to target nodes when the targets can be resolved unambiguously.
In parallel, it records certificate-field concepts and aliases, namely subjectAltName, dNSName, basicConstraints.cA, and policy OIDs, to map natural-language mentions to canonical certificate paths during later extraction.

PKI specifications form a densely interconnected cross-document network.
For example, RFC~5280 references multiple RFCs, while CABF~BR directly depends on RFC~5280 field definitions.
We model this knowledge as a multi-edge directed graph with eight node types, namely Specification, Section, Definition, CertificateField, Rule, Operation, Value, and Concept, and four source-backed relations, namely CONTAINS, DEFINES, REFERENCES, and APPLIES\_TO. Appendix~\ref{app:kg} summarizes these relations and their deterministic retrieval constraints.

This builder is deliberately conservative.
It stores only relations directly supported by document structure or textual evidence, while leaving normative judgments such as conflict, override, and lintability to later deterministic stages.
Normative-judgment relations such as CONFLICTS\_WITH and OVERRIDES are produced independently by the rule engine and are not included in the retrieval graph.
The resulting graph serves as a traceable context index rather than as an oracle. It helps the extractor identify definitions, inherited constraints, and referenced rules, but does not determine whether a requirement is lintable.

\subsection{Deterministic Context Retrieval}
\label{sec:method-graphrag}

To address cross-document references, DLEP organizes context deterministically rather than relying on the LLM to infer relations.
We implement a GraphRAG~\cite{edge2024local} retrieval module that treats the knowledge graph of \S\ref{sec:method-kg-builder} as a fixed input and retrieves context through bounded subgraph queries.

Given a target Section node, the retrieval module performs a BFS expansion over the $k$-hop neighborhood using the source-backed CONTAINS, DEFINES, REFERENCES, and APPLIES\_TO relations.
It then assembles LLM context by relation-type priority and provides term definitions, field metadata, and referenced sections. The retrieval process uses neither vector similarity nor LLM inference, and produces no independent judgment.

\subsection{DLEP}
\label{sec:method-dual-layer}

\begin{table}[tb]
\centering
\caption{RFC~2119 keyword tiers.}
\label{tab:rfc2119}
\setlength{\tabcolsep}{4pt}
\begin{tabular}{@{}L{0.27\linewidth} L{0.30\linewidth} L{0.33\linewidth}@{}}
\hline
Tier & Keywords & Semantics \\
\hline
Mandatory      & MUST, SHALL, REQUIRED   & Absolute requirement; a violation is non-compliance \\
Prohibition    & MUST NOT, SHALL NOT     & Absolute prohibition \\
Recommendation & SHOULD, SHOULD NOT, RECOMMENDED & Strongly recommended or discouraged; justified deviation permitted \\
Optional       & MAY, OPTIONAL           & Fully optional \\
\hline
\end{tabular}
\end{table}

\begin{table}[tb]
\centering
\caption{IR core four-tuple fields.}
\label{tab:ir-core}
\setlength{\tabcolsep}{4pt}
\footnotesize
\begin{tabular}{@{}L{0.18\linewidth} L{0.74\linewidth}@{}}
\hline
Field & Meaning \\
\hline
subject    & Certificate field path (e.g., extensions.\allowbreak subjectAltName.\allowbreak dNSName), resolved by FieldResolver \\
obligation & RFC~2119 keyword (MUST/SHALL/SHOULD/\dots) \\
predicate  & Assertion type (must\_be\_present, conform\_to, equal, \dots) \\
constraint & constraint value or pattern \\
\hline
\end{tabular}
\end{table}

To address the LLM controllability challenge, we design a dual-layer extraction pipeline (DLEP) that preserves the LLM's semantic understanding while keeping the extraction process auditable.

Layer~1 performs deterministic recall.
Web PKI specifications commonly express normative requirements using the RFC~2119 keywords listed in Table~\ref{tab:rfc2119}.
RFC~8174 further specifies that these keywords carry normative force only when capitalized, enabling deterministic keyword matching.
Since ETSI~EN~319 does not follow this convention, its candidates are recalled through lowercase matching.

To maximize recall, Layer~1 operates in three passes.
Pass~1 directly matches RFC~2119 keywords.
Pass~2 detects nested structures so subordinate rules inherit the parent obligation level.
Pass~3 captures normative statements that do not follow the standard RFC~2119 form.

Layer~2 performs controlled semantic parsing using an LLM, but restricts the model to a constrained language-understanding role.
Instead of generating lint rules or compliance judgments directly, Layer~2 converts normative text into a JSON-serialized intermediate representation (IR), while all normative reasoning is deferred to later deterministic stages.

The IR decouples semantic extraction from downstream tasks, allowing the same extraction result to support multiple uses, as shown in Fig.~\ref{fig:ir-downstream}.
Each IR record is serialized as a JSON object with 33 fields.
Its core component, shown in Table~\ref{tab:ir-core}, is the four-tuple
$\mathrm{IR} = \langle\text{subject}, \text{obligation}, \text{predicate}, \text{constraint}\rangle$,
together with 29 extension fields summarized in Appendix~\ref{app:ir}.

This IR-based design has three advantages over direct end-to-end generation of lint checks from raw text.
First, citation, constraint, and obligation information are stored separately, improving traceability and conflict handling.
Second, the IR serves as deterministic input to the four-condition lintability analysis in \S\ref{sec:method-lintability}, making lintability decisions reproducible rather than model-dependent.
Third, once a rule is determined to be lintable, generating static lint checks from the IR is more stable and interpretable because key elements such as field paths, assertion types, and constraint values are already explicitly structured.

To constrain LLM behavior, Layer~2 enforces a predefined schema and category set.
Every extracted rule must map to explicit source text, and certificate field hierarchies must follow the ASN.1 path definitions in RFC~5280.
The IR generation process is also staged. The model first determines the rule category and then fills in the remaining fields, which reduces errors compared with generating all fields in a single step.
Appendix~\ref{app:ir} provides the category definitions, correction rules, and tuning details.

\begin{figure}[tb]
  \centering
  \cicasfig{figures/fig02_ir_downstream.png}{IR-driven ``extract once, apply
    many'' architecture.}
  \caption{IR-driven ``extract once, apply many'' architecture.}
  \label{fig:ir-downstream}
\end{figure}

\subsection{Normative Rule Lintability Decision}
\label{sec:method-lintability}
\label{sec:lintability}

\begin{table*}[tb]
\centering
\caption{The four lintability conditions, each a Boolean function of one IR field.}
\label{tab:four-conditions}
\footnotesize
\setlength{\tabcolsep}{4pt}
\renewcommand{\arraystretch}{1.2}
\begin{tabular}{@{}l l L{0.36\linewidth} L{0.30\linewidth}@{}}
\hline
ID & IR field & Value domain & Condition \\
\hline
$C_1$      & obligation         & MUST, SHALL, REQUIRED, SHOULD, RECOMMENDED, MUST NOT, SHALL NOT, SHOULD NOT, MAY, OPTIONAL & obligation $\notin$ \{MAY, OPTIONAL\} \\
$C_2$      & assertion\_subject & Certificate, CA, CrossArtifact & assertion\_subject $=$ Certificate \\
$C_3$      & enforcement\_phase & Encoding, Validation, Comparison, Processing & enforcement\_phase $=$ Encoding \\
$\neg C_4$ & rule\_category     & encoding\_constraint, comparison, definition, capability, algorithm\_ref, display, \dots & rule\_category $\notin$ \{definition, capability, algorithm\_ref, display, \dots\} \\
\hline
\end{tabular}
\end{table*}

Before generating code, we decide the lintability of each normative rule, namely whether it admits such a static lint check. Rather than asking the LLM to make this judgment directly, we compute it from the IR that the LLM has already produced. Three observations motivate this design. First, the same four pieces of information determine the lintability of normative rules in practice, namely the deontic strength of the obligation, who is obligated, when the obligation applies, and what type of rule it is. Second, each of the four can be encoded as a single IR field with a small, closed value set, so the lintability decision becomes a Boolean function of four discrete fields. Third, separating the four conditions makes any wrong label traceable to exactly one IR field rather than to an opaque end-to-end LLM call.

Table~\ref{tab:four-conditions} lists the four IR fields, their value domains, and the conditions they decide. The deontic level obligation comes from the IR core four-tuple in Table~\ref{tab:ir-core}. The remaining three fields, assertion\_\allowbreak subject, enforcement\_\allowbreak phase, and rule\_\allowbreak category, are key extension fields. From the enumeration of rule\_\allowbreak category, we designate
\begin{gather*}
\mathcal{N} = \{definition,\,capability,\\
algorithm\_ref,\,display,\,\dots\}
\end{gather*}
as the subset whose rule type carries no static certificate lint check on the certificate. This subset covers term definitions, CA capability statements, delegations to external algorithm specifications, UI rendering rules, and similar non-encoding rule types.

The four conditions in Table~\ref{tab:four-conditions} are necessary for lintability, and each is a Boolean function of exactly one IR field:
\begin{equation}\label{eq:lintable}
\mathrm{lintable}(r) \;\Longrightarrow\; C_1(r)\wedge C_2(r)\wedge C_3(r)\wedge \neg C_4(r),
\end{equation}
\begin{align*}
C_1(r) &\equiv \mathrm{is\_normative}(r.\text{obligation}), \\
C_2(r) &\equiv r.\text{assertion\_subject} = Certificate, \\
C_3(r) &\equiv r.\text{enforcement\_phase} = Encoding, \\
C_4(r) &\equiv r.\text{rule\_category} \in \mathcal{N}.
\end{align*}
The obligation field also fixes the lint severity. MUST-class obligations map to Error, and SHOULD-class obligations map to Warning. Severity is therefore not a second classifier, but a direct read of obligation.

The three non-deontic conditions encode three orthogonal boundaries. $C_2$ is the subject boundary, separating obligations on the certificate from obligations on the CA, relying party, or surrounding ecosystem. $C_3$ is the runtime boundary, separating obligations observable in the issued bytes from those that fire during chain processing, name comparison, revocation handling, or CAA retrieval. $\neg C_4$ is the process boundary, separating state assertions from terms in $\mathcal{N}$.

Since each $C_i$ is a deterministic function of one IR field, the lintability
decision---which adopts the conjunction $C_1\wedge C_2\wedge C_3\wedge\neg C_4$ of
these necessary conditions as its operational criterion---is computed from
the IR rather than predicted by the LLM. The same normative rule IR yields the same label bit-for-bit on every run, and any misclassification is traceable to a specific Layer~2 field assignment.

%-------------------------------------------------------------------------------
\section{Restricted Code Generation}
\label{sec:codegen}
%-------------------------------------------------------------------------------

Generating checking code directly with a large language model faces two
inherent difficulties. The first is hallucination. The model may
misspell a field name as another field that happens to exist but is
semantically unrelated, or introduce a unit error in an OID. The second is
output instability. The same code-generation prompt often produces
different outputs across calls, with no guarantee that the result will compile or
be logically correct. To contain this uncertainty, we confine a rule's code body
to a finite, closed DSL-tree space $\Tv$. Each rule's code body can only be a
tree in that space, which is then rendered into Go code by a deterministic
function. The deterministic main path avoids unconstrained code emission, and
the restricted LLM fallback is bounded to the same closed space and gated by
``render and pass go build,'' so any out-of-space or non-compiling output is
rejected rather than shipped.
Figure~\ref{fig:codegen-pipeline} illustrates the whole
pipeline on a single normative rule---CA certificates must assert the keyCertSign
key-usage bit.

\begin{figure*}[tb]
  \centering
  \paperfig{figures/fig03_codegen_pipeline.png}{Worked example of restricted
    code generation. A single normative rule flows from its filled IR fields through the
    DSL tree, represented as an abstract syntax tree, and its JSON serialization to deterministically rendered Go
    checking code.}
  \caption{Worked example of restricted code generation. }
  \label{fig:codegen-pipeline}
\end{figure*}

\subsection{Atomic Templates and IR Filling}
\label{subsec:atoms}

An atomic template is the smallest unit of a certificate lint.
Most certificate lints in PKI are logical combinations of a
few atomic templates that test whether an extension is present, whether a
field equals a constant, or whether every element of a list matches a
pattern, and combining such atomic templates with conjunction, disjunction, and
negation can
characterize most static constraints. The abstract syntax of the code DSL is:
\begin{equation}\label{eq:grammar}
  \mathcal{T} \;::=\; a(\bar{v}) \;\mid\; \neg\, \mathcal{T} \;\mid\;
  \mathcal{T} \wedge \mathcal{T} \;\mid\; \mathcal{T} \vee \mathcal{T}
\end{equation}
where $a \in \mathcal{A}$ is an atomic-template predicate and $\bar{v}$ is its
argument list. $\mathcal{A}$ is a finite, closed set of atomic
templates; the experiments use 151 atomic templates in total, split by
generality into 98 GENERIC and 53 NON\_GENERIC, with grading criteria and
representative examples in Appendix~\ref{app:tv}; $\{\neg, \wedge, \vee\}$ are
the propositional-logic combinators. The code body of a lint rule is modeled as
an ordered pair $(p, q) \in \mathcal{T}_\varepsilon \times \mathcal{T}$, where
$q \in \mathcal{T}$ is the main assertion and $p \in \mathcal{T}_\varepsilon =
\mathcal{T} \cup \{\varepsilon\}$ is an optional precondition, where $\varepsilon$
marks an absent precondition. In the
equation below, $c$ is a single certificate, and $\lVert u \rVert(c) \in
\{\mathrm{true}, \mathrm{false}\}$ denotes the Boolean result of evaluating a
DSL tree or precondition $u$ on $c$; atomic templates evaluate by their semantics, and
$\neg/\wedge/\vee$ by classical propositional logic. The execution semantics
is:
\begin{equation}\label{eq:exec}
  \lVert (p, q) \rVert(c) \;=\;
  \begin{cases}
    \mathrm{NA}, & p \neq \varepsilon \,\land\, \lnot\lVert p \rVert(c) \\[2pt]
    \mathrm{Pass}, & (p = \varepsilon \,\lor\, \lVert p \rVert(c)) \,\land\,
      \lVert q \rVert(c) \\[2pt]
    \mathrm{Severity}(r), & \text{otherwise}
  \end{cases}
\end{equation}
The precondition $p$ is the normative rule's precondition. When it exists and is
false on $c$, the result is $\mathrm{NA}$, not applicable; when the
precondition is absent or holds and $q$ holds, the result is $\mathrm{Pass}$;
when the precondition holds but $q$ does not, it is a violation, returning
$\mathrm{Severity}(r)$ by obligation. Because $C_1$ of \S\ref{sec:lintability}
already excludes MAY/OPTIONAL, an obligation entering generation must be
MUST-class or SHOULD-class, so $\mathrm{Severity}(r) \in \{\mathrm{Error},
\mathrm{Warn}\}$.

\textbf{The IR fills the atomic template's parameter slots.} Selecting an atomic
template fixes only which check to use; one still needs to know what to check and
against what value, and these parameter values come precisely from the current
normative rule's IR. The IR's subject, the certificate field path given by the
field resolver such as subjectAltName.dNSName, fills the atomic
template's field slot, and the IR's constraint, a constraint value or
pattern such as a value set, a length interval, or a regex name, fills the value
slot. For example, the normative rule ``the dNSName of subjectAltName must be non-empty''
is, after mapping and filling, instantiated as
$\mathrm{FieldNonEmpty}(\mathrm{DNSNames})$.

\textbf{Mapping IR predicates to atomic templates.} The upstream
IR and the downstream DSL use two vocabularies. The IR describes normative rules
with extraction-stage predicates such as must\_be\_present,
encode\_as, and in\_range, while the DSL expresses checks with
atomic templates, and the two are not in one-to-one correspondence. To
bridge them, we maintain a many-to-many predicate-to-template map that assigns each IR
predicate a set of candidate atomic templates that can semantically carry it.
For instance, must\_be\_present may correspond to $\{\mathrm{ExtPresent},
\mathrm{FieldNonEmpty}\}$, and in\_range to $\{\mathrm{IntInRange},
\mathrm{PathLenConstraintHas}, \dots\}$---which one to take depends on the type
of the constrained field.

\subsection{Typed Vocabulary and Parameter Closure}
\label{subsec:tv}

\S\ref{subsec:atoms} described how atomic templates and IR filling place the
IR's subject and constraint into an atomic template's parameter slots---but the
IR provides only the field that a normative rule restricts and the value it
restricts that field to,
without constraining the legal type of those values. If the subject path
in the IR points to a nonexistent field, or a string is placed into a slot that
requires an integer, a structurally legal but semantically misaligned tree can
still result. This subsection closes that gap with a typed vocabulary
$\mathcal{V}$. $\mathcal{V}$ is the disjoint union of seven finite sets, frozen
at system startup, and each atomic-template parameter may only draw from one of
them---certificate field names, DN fields, OID constants, KeyUsage bits,
ExtKeyUsage bits, ASN.1 encoding types, and named regexes. The size and samples
of each class appear in Appendix~\ref{app:tv}.

Each atomic template has a signature table specifying which class, or base type
integer/Boolean/string, each of its parameters must draw from; a call to an
atomic template is legal if and only if every actual argument falls in the class its
signature prescribes. Writing $\Tv$ for the set of DSL trees all of whose
parameters are legal, we strictly restrict the range of the code generator to
it. Any tree in $\Tv$ contains no field name, OID, or regex
outside $\mathcal{V}$; an out-of-bounds tree necessarily errors at the parsing
stage, triggering repair, and never enters code generation.

\subsection{Code Generation and Mechanical Translation}
\label{subsec:mech}

A legal DSL tree $t$ must be turned into two products. One is executable Go
checking code, and the other is a code summary, the code\_summary, that
summarizes its checking logic in one sentence. Both conversions are performed by
deterministic total functions, with no LLM call. The rendering function $\rho$
maps a tree to Go code, and the mechanical translation function $\smech$ maps the
same tree to a code\_summary sentence; for the same $t$, each produces a unique
output.

$\rho$ is type-safe. It never generates an out-of-bounds field, OID, or
parameter type, so closure is preserved from the DSL-tree layer down to the
Go-code layer. Together with determinism, this lets the certificate-level
execution verification of \S\ref{subsec:exec-verify} treat $\rho(t)$ as a fixed
function of $t$, whose execution behavior does not vary with run count or model.
$\smech$, by contrast, is structure-preserving. Each atomic template and
combinator maps to a distinct phrase, so $\smech(t)$ retains the structural
information of the DSL tree in the verification chain; the construction is
detailed in Appendix~\ref{app:sigma}.

Thus $\rho$'s product leads to executable code and the certificate-level execution
verification of \S\ref{subsec:exec-verify}, while $\smech$'s product leads to
semantic-equivalence verification---comparing the code\_summary with the
normative rule text and determining whether the two are semantically equivalent.

\subsection{Tree Synthesis with a Deterministic Main Path and Restricted LLM
Fallback}
\label{subsec:synth}

The preceding discussion described how the IR filling of \S\ref{subsec:atoms} and the
parameter closure of \S\ref{subsec:tv} confine a DSL tree to $\Tv$, but two
decisions still require synthesis. The first is which atomic template carries
which IR predicate. The second is whether the IR's subject and
constraint can both be filled into the chosen atomic template.
This subsection gives the implementation of the generator. It
first attempts deterministic synthesis and falls back to LLM synthesis only when
deterministic synthesis returns no tree. Both paths use rendering and passing go build as
the acceptance gate---even if deterministic synthesis returns a tree, that is
not enough; it must actually compile, otherwise it is treated as unresolved and
passed to LLM synthesis. Both paths are constrained by the same set of
structural checks, and neither determines whether the code expresses the
normative rule, which is the job of the verification chain of
\S\ref{subsec:exec-verify}; they only confine the output to a traceable space.

\textbf{Deterministic main path.} When an IR predicate selects an atomic
template through the predicate-to-template map and its subject and constraint can fill that template's
parameter slots, the generator deterministically derives a DSL tree using
predefined rules. Here, deriving means constructing a DSL tree from the IR step by
step by fixed rules, with no LLM call throughout. Most lintable normative rules in our
study are covered by this deterministic path.

\textbf{Restricted LLM synthesis.} When deterministic synthesis returns no tree,
the generator switches to LLM synthesis. The model receives the normative-rule context,
the structured IR, the candidate atomic-template set, and a readable enumeration
and signature table of $\mathcal{V}$ and $\mathcal{A}$. It
outputs only one of two results, either a serialized DSL tree or an explicit
no-template abstention marker NT. The raw LLM
output is first validated by a parser, which resolves it into
a legal DSL tree $t \in \Tv$, a type error, or NT, and any
identifier outside $\mathcal{V}/\mathcal{A}$ is rejected here; a passing tree is
rendered by $\rho$, and then a deterministic post-processor
binds Description/Citation/Name and the metadata of
each normative source, with details in Appendix~\ref{app:tv}.

\subsection{Three-Layer Semantic Alignment Verification}
\label{subsec:exec-verify}

Whether the generated code expresses the normative rule is the hardest judgment
in the whole pipeline. This subsection checks it at three levels from coarse to
fine.

\textbf{Layer~A, description provenance.} Check whether the Description
can be traced back to the normative source.

\textbf{Layer~B, semantic equivalence of the code-behavior summary.} It turns
the cross-modal judgment of whether the code implements the
normative rule into the same-modal judgment of whether two natural-language
sentences are semantically equivalent. Following $\smech$ of \S\ref{subsec:mech}, with the
phrase dictionary in Appendix~\ref{app:sigma}, the code is translated into a
code\_summary, which is then judged against the Description.

\textbf{Layer~C, compilation and structure.} Verify that the code can be parsed
syntactically, that the function and metadata fields are complete, that
dependencies are correctly imported, and that the lint is registered.

The three layers form a Boolean acceptance gate by conjunction. Layer~B's
judgment is the core signal, and Layer~A's provenance serves as a
prefilter; failing any layer triggers the repair loop of Appendix~\ref{sec:saiv}.
Since $\rho$ and $\smech$ are deterministic and Layer~A anchors the Description
to source text, the LLM contributes only this one same-modal comparison, and an
IR error surfaces here as a mismatch.

\textbf{Certificate-level execution verification, execution-level evidence for
\CodeEqIR.} Layer~B's semantic-equivalence judgment is given by an LLM and cannot be fully
determinized; as a complement, we introduce a model-independent
deterministic verification to check whether the generated code executes the DSL
tree it corresponds to, written \CodeEqIR. Its idea borrows
from the test oracle in software testing~\cite{barr2015oracle}, using a set of samples with
known-correct answers to check the code under test. It proceeds in two steps.

Step 1, checking a single atomic template. For each atomic template we
construct a pair of controlled certificate samples, one that should
make the template's predicate true, and one that should make it false. The
samples are constructed by the atomic template's parameter types and are not
bound to any concrete normative rule. The rendered code is executed for real on the two
certificates and the verdicts are read back; the atomic template is accepted by
this check if and only if the results on both certificates match expectation.

Step 2, generalizing from atomic templates to the whole tree. For a
generated lint holding a DSL tree $t$, if $t$ is composed only of accepted
atomic templates and compiles as a whole, we can induct over the tree
structure. The behavior of the leaves is checked in Step~1, and the combinators
$\neg/\wedge/\vee$ evaluate by classical propositional logic and preserve that
behavior level by level, so the rendered code $\rho(t)$ executes the whole tree,
written $\mathrm{Code}\equiv t$; and because $t$ is deterministically reduced from
the IR, \CodeEqIR.

This verification is performed on controlled certificate samples; the result is
deterministic and does not vary with run count or model. Its boundary
must be stressed. \CodeEqIR{} only shows that the code follows the intermediate
representation IR, which is not the same as code--normative-rule
alignment; if the IR itself has already drifted from the normative rule, this
verification cannot detect it. The code--normative-rule judgment therefore still
rests with \CodeEqRule{}, the Layer~B semantic-equivalence judgment.

%-------------------------------------------------------------------------------
\section{Experimental Evaluation}
\label{sec:eval}
%-------------------------------------------------------------------------------

This section runs experiments with the methods and framework above and draws
quantitative conclusions. Besides the internal layered metrics
(the end-to-end result snapshot below and the lint-coverage analysis of
\S\ref{subsec:coverage}), it provides two
mutually independent external validations, each checking one half of the
framework: \S\ref{subsec:lintability-ext} externally validates rule
extraction and lintability determination against the zlint maintainers' manual
gold standard, and \S\ref{subsec:certdetect} externally validates code
generation and synonymy judgment by execution on real certificates.

\subsection*{Experimental Setup and Result Snapshot}

Evaluation is performed on full-corpus re-extraction results of the two standard
sources RFC~5280 and CABF BR. The LLM used in the experiments is GPT-5.4.

\begin{table*}[tb]
\centering
\footnotesize
\caption{Current end-to-end result snapshot (RFC~5280 $+$ CABF BR).}
\label{tab:snapshot}
\begin{tabular}{@{}p{4.4cm}r p{10.4cm}@{}}
\toprule
Metric & Count & Scope / notes \\
\midrule
Total recalled rules & 2080 & RFC~5280 640 $+$ CABF BR 1440 \\
lintable (single-cert observable) & 248 &
  current re-extraction / re-adjudication: CABF 170 $+$ RFC~5280 78; the negative
  gate excludes process, cross-certificate, runtime, external-semantic, and
  MAY/OPTIONAL classes; CRL document rules are outside this scope \\
zlint coverage (full; existing cognate lint) & 158 &
  the rule is fully implemented by some cognate zlint certificate lint; residual
  cross-standard inheritance gaps are analyzed in \S\ref{subsec:crossstd} \\
uncovered lintable (codegen domain) & 90 &
  enters the current codegen domain ($248 - 158$); CABF 64 $+$ RFC~5280 26, the
  target set of downstream code generation $\phi_G$ \\
can generate compilable lint & 90 &
  deterministic path 87 $+$ LLM fallback path 3, both gated by ``render and
  pass go build''; generation rate $= 90/90 = 100.0\%$ \\
cannot generate (no tree / not compiling / abstain) & 0 &
  no generation failure in the current codegen domain \\
final-shipping strict synonymous (\CodeEqSpec) & 90 &
  the final in-tree zlint lint (\texttt{CheckApplies} $+$ \texttt{Execute} $+$
  severity/date metadata) judged against the original rule text and context;
  strict synonymy rate $= 90/90 = 100.0\%$ \\
final-shipping not synonymous (does not express) & 0 &
  no final-shipping residual in the current codegen domain \\
emitted lints' atom templates touched & 52 &
  28 GENERIC $+$ 24 NON\_GENERIC of the 136 registered atoms; by rule, 53
  GENERIC-only / 5 mixed / 32 NON\_GENERIC-only \\
\emph{(collateral) certificate-level execution verification \CodeEqIR} &
  \emph{37} &
  \emph{model-independent IR-level faithfulness, not a synonymy numerator; the
  paper-facing synonymy rate uses only the final emitted code versus the source
  rule} \\
\bottomrule
\end{tabular}
\end{table*}

Under the current codegen domain there is no non-generable compilable lint and no
final-shipping non-synonymous residual (both counts are 0). The 37 rules
certified \CodeEqIR{} by certificate-level execution verification are reported
only as collateral IR-level evidence: \CodeEqIR{} is model-independent and
deterministic, yet it does not by itself imply \CodeEqSpec, so the paper-facing
synonymy rate is taken solely from the final emitted code judged against the
source rule. Any reduction of the denominator therefore comes from re-extraction
and lintability re-adjudication rather than from local metric filters; e.g.,
R31068 (RFC~5280~\S4.2.1.6) was re-adjudicated as not lintable because a single
certificate exposes only the chosen GeneralName tag and its value, not the
issuer's pre-encoding intent.

The first-layer quality signal is G1 conservation (Theorem~\ref{thm:recall},
\S\ref{subsec:targets}): the total number of keyword-recalled rules should
neither increase nor decrease in any downstream classification step. Under the
current snapshot, conservation holds strictly:
\[
  \underbrace{2080}_{\text{recall}} \;=\; \underbrace{529}_{\text{noise}} +
  \underbrace{1551}_{\text{genuine}}, \quad
  \underbrace{1551}_{\text{genuine}} \;=\; \underbrace{248}_{\text{lintable}} +
  \underbrace{1303}_{\text{not lintable}}.
\]
Both identities hold term by term, i.e., the G1 residual $\Lrec = 0$ and the top
two layers of conservation are closed (deterministically recomputable). The 248
lintable rules split by standard source
into 170 CABF and 78 RFC~5280; of these zlint fully covers 158, leaving
90---the domain of downstream code generation $\phi_G$.

\subsection{Lint Coverage Analysis}
\label{subsec:coverage}

Answering ``how many lintable rules do existing tools cover'' hinges on using
the right criterion. We use ``whether the lintable rule is genuinely implemented
by some cognate zlint lint'' as the criterion: for each rule we retrieve
candidate zlint lints, then compare field by field (subject / obligation /
predicate / constraint), giving a three-level verdict full (fully implemented) /
partial (partially implemented) / none (not implemented).

Fig.~\ref{alg:coverage} formalizes this determination: summarize a lint
into a reverse IR, then align field by field; candidate retrieval is by
source/section (Stage-1: RFC section numbers are stable and narrowed by prefix,
while CABF section numbers drift across versions so all CABF lints are taken),
with an added ``wrong-field'' consistency gate (Stage-3, only downgrades), and
coverage counts only full (Stage-4; partial/none go to $\phi_G$).

% Algorithm 1 -- zlint coverage-alignment determination for lintable normative rules.
% Included from sections/08_evaluation.tex via % Algorithm 1 -- zlint coverage-alignment determination for lintable normative rules.
% Included from sections/08_evaluation.tex via \input{algorithms/alg_coverage}.
%
% Typeset single-column (not figure*) at \scriptsize: the algorithmic lines are
% short, so at full width the right half of most lines was blank and the float was
% billed two columns of page area. Long \Comment texts are placed on their own
% \Statex lines because \Comment is right-aligned and cannot wrap. No step, symbol,
% or comment wording was changed.
\begin{figure}[tb]
\makeatletter
\@ifundefined{c@algorithm}{%
  \newcounter{algorithm}%
  \renewcommand{\thealgorithm}{\arabic{algorithm}}%
}{}%
\makeatother
\refstepcounter{algorithm}
\label{alg:coverage}
\noindent\rule{\linewidth}{0.6pt}
\par\smallskip
\noindent\textbf{Algorithm~\thealgorithm.} Algorithmic listing for zlint
coverage-alignment determination for individual normative rules.
\par\smallskip
\noindent\rule{\linewidth}{0.4pt}
\par\smallskip
\scriptsize
\begin{algorithmic}[1]
\Statex \textbf{Input:} lintable normative-rule set $B=\{b_1,\dots,b_m\}$, where each normative rule has
  IR four-tuple $\langle$subject, obligation, predicate, constraint$\rangle$,
  source, and section; zlint lint set $L=\{\ell_1,\dots,\ell_n\}$, where each lint has
  description/citation/severity metadata and Pass/Error tests; offline code
  summarizer $M_{\mathrm{summ}}$, and field-level judge $M_{\mathrm{judge}}$
  with temperature $=0$
\Statex \textbf{Output:} per-normative-rule coverage verdict $\mathrm{verdict}(b)\in\{\text{full, partial,
  none}\}$, covered set $\RLcov$, and uncovered code-generation domain $\RLunc$
\Statex // Offline cached Stage-0: summarize each lint into a reverse
  IR
\ForAll{$\ell \in L$}
  \State $\ell.\mathit{ir} \gets M_{\mathrm{summ}}(\ell.\mathrm{src},
    \ell.\mathrm{tests}, \ell.\mathrm{meta})$
  \Statex \Comment{four-tuple $+$ code\_summary}
\EndFor
\State pre-split candidate pools by Source: $L_{\mathrm{RFC}}\gets$ RFC lints,\;
  $L_{\mathrm{CABF}}\gets$ CABF lints
\ForAll{$b \in B$}
  \State // Stage-1: deterministic candidate retrieval
  \If{$b.\mathrm{source}=\mathrm{RFC}$}
    \State $C(b)\gets\{\ell\in L_{\mathrm{RFC}}:
      \textsc{SectionPrefixMatch}(b.\mathrm{section},\ell.\mathrm{citedSections})\}$
  \ElsIf{$b.\mathrm{source}=\mathrm{CABF}$}
    \State $C(b)\gets L_{\mathrm{CABF}}$ \Comment{CABF sections drift $\Rightarrow$ take all}
  \Else
    \State $C(b)\gets\emptyset$
  \EndIf
  \State // Stage-2: field-level judge, compare four-tuple, best across candidates
  \State $(v,\ell^{*})\gets M_{\mathrm{judge}}(b.\mathit{ir}, C(b))$
  \Statex \Comment{$v\in\{\text{full},\text{partial},\text{none}\}$}
  \State // Stage-3: reject positive matches with conflicting concrete subject families
  \If{$v\in\{\text{full},\text{partial}\} \wedge
      \mathrm{Family}(b.\mathrm{subject})\neq\emptyset \wedge
      \mathrm{Family}(\ell^{*}.\mathrm{subject})\neq\emptyset \wedge
      \mathrm{Family}(b.\mathrm{subject})\neq\mathrm{Family}(\ell^{*}.\mathrm{subject})$}
    \State $v\gets\text{none}$
    \Statex \Comment{different concrete families $\Rightarrow$ wrong-field false positive}
  \EndIf
  \State $\mathrm{verdict}(b)\gets v$
\EndFor
\State // Stage-4: deterministic aggregation, coverage counts full only
\State $\RLcov\gets\{b\in B:\mathrm{verdict}(b)=\text{full}\}$;\quad
  $\RLunc\gets B\setminus \RLcov$
\State \textbf{return} $(\{(b,\mathrm{verdict}(b)):b\in B\},\ \RLcov,\ \RLunc)$
\end{algorithmic}
\par\smallskip
\noindent\rule{\linewidth}{0.6pt}
\end{figure}
.
%
% Typeset single-column (not figure*) at \scriptsize: the algorithmic lines are
% short, so at full width the right half of most lines was blank and the float was
% billed two columns of page area. Long \Comment texts are placed on their own
% \Statex lines because \Comment is right-aligned and cannot wrap. No step, symbol,
% or comment wording was changed.
\begin{figure}[tb]
\makeatletter
\@ifundefined{c@algorithm}{%
  \newcounter{algorithm}%
  \renewcommand{\thealgorithm}{\arabic{algorithm}}%
}{}%
\makeatother
\refstepcounter{algorithm}
\label{alg:coverage}
\noindent\rule{\linewidth}{0.6pt}
\par\smallskip
\noindent\textbf{Algorithm~\thealgorithm.} Algorithmic listing for zlint
coverage-alignment determination for individual normative rules.
\par\smallskip
\noindent\rule{\linewidth}{0.4pt}
\par\smallskip
\scriptsize
\begin{algorithmic}[1]
\Statex \textbf{Input:} lintable normative-rule set $B=\{b_1,\dots,b_m\}$, where each normative rule has
  IR four-tuple $\langle$subject, obligation, predicate, constraint$\rangle$,
  source, and section; zlint lint set $L=\{\ell_1,\dots,\ell_n\}$, where each lint has
  description/citation/severity metadata and Pass/Error tests; offline code
  summarizer $M_{\mathrm{summ}}$, and field-level judge $M_{\mathrm{judge}}$
  with temperature $=0$
\Statex \textbf{Output:} per-normative-rule coverage verdict $\mathrm{verdict}(b)\in\{\text{full, partial,
  none}\}$, covered set $\RLcov$, and uncovered code-generation domain $\RLunc$
\Statex // Offline cached Stage-0: summarize each lint into a reverse
  IR
\ForAll{$\ell \in L$}
  \State $\ell.\mathit{ir} \gets M_{\mathrm{summ}}(\ell.\mathrm{src},
    \ell.\mathrm{tests}, \ell.\mathrm{meta})$
  \Statex \Comment{four-tuple $+$ code\_summary}
\EndFor
\State pre-split candidate pools by Source: $L_{\mathrm{RFC}}\gets$ RFC lints,\;
  $L_{\mathrm{CABF}}\gets$ CABF lints
\ForAll{$b \in B$}
  \State // Stage-1: deterministic candidate retrieval
  \If{$b.\mathrm{source}=\mathrm{RFC}$}
    \State $C(b)\gets\{\ell\in L_{\mathrm{RFC}}:
      \textsc{SectionPrefixMatch}(b.\mathrm{section},\ell.\mathrm{citedSections})\}$
  \ElsIf{$b.\mathrm{source}=\mathrm{CABF}$}
    \State $C(b)\gets L_{\mathrm{CABF}}$ \Comment{CABF sections drift $\Rightarrow$ take all}
  \Else
    \State $C(b)\gets\emptyset$
  \EndIf
  \State // Stage-2: field-level judge, compare four-tuple, best across candidates
  \State $(v,\ell^{*})\gets M_{\mathrm{judge}}(b.\mathit{ir}, C(b))$
  \Statex \Comment{$v\in\{\text{full},\text{partial},\text{none}\}$}
  \State // Stage-3: reject positive matches with conflicting concrete subject families
  \If{$v\in\{\text{full},\text{partial}\} \wedge
      \mathrm{Family}(b.\mathrm{subject})\neq\emptyset \wedge
      \mathrm{Family}(\ell^{*}.\mathrm{subject})\neq\emptyset \wedge
      \mathrm{Family}(b.\mathrm{subject})\neq\mathrm{Family}(\ell^{*}.\mathrm{subject})$}
    \State $v\gets\text{none}$
    \Statex \Comment{different concrete families $\Rightarrow$ wrong-field false positive}
  \EndIf
  \State $\mathrm{verdict}(b)\gets v$
\EndFor
\State // Stage-4: deterministic aggregation, coverage counts full only
\State $\RLcov\gets\{b\in B:\mathrm{verdict}(b)=\text{full}\}$;\quad
  $\RLunc\gets B\setminus \RLcov$
\State \textbf{return} $(\{(b,\mathrm{verdict}(b)):b\in B\},\ \RLcov,\ \RLunc)$
\end{algorithmic}
\par\smallskip
\noindent\rule{\linewidth}{0.6pt}
\end{figure}


Here \textsc{SectionPrefixMatch} is true when the rule's section number shares a
common prefix with a lint's cited section number; $\mathrm{Family}(\cdot)$ maps a
subject path to a concrete field family, returning $\emptyset$ for a
vague/unresolved subject rather than downgrading (so genuine matches on coarse
subjects are preserved). The algorithm calls the LLM only in Stage-0 and Stage-2
(the judge at temperature $=0$, its prompt containing three kinds of
positive/negative examples: direction reversal, field misplacement, and
constraint-type confusion); summaries are cached offline, the full CABF
candidate set is sent to the judge in batches, and the rest is deterministic.

Of the 248 lintable rules, zlint fully covers 158 and leaves 90
uncovered under this coverage procedure (Table~\ref{tab:coverage}); the 90
uncovered are the domain of code generation $\phi_G$.

\begin{table}[tb]
\centering
\small
\caption{zlint cognate lint counts and their coverage of the 248 lintable
rules, by standard source.}
\label{tab:coverage}
\begin{tabular}{@{}lrrr@{}}
\toprule
Item & CABF & RFC~5280 & Total \\
\midrule
\emph{zlint cognate lints, total (ref.)} & \emph{172} & \emph{123} & \emph{295} \\
\quad\emph{— of which cert.\ lints} & \emph{166} & \emph{115} & \emph{281} \\
\quad\emph{— of which CRL lints} & \emph{6} & \emph{8} & \emph{14} \\
full (full coverage) & 106 & 52 & 158 \\
uncovered (codegen domain) & 64 & 26 & 90 \\
lintable total & 170 & 78 & 248 \\
\bottomrule
\end{tabular}
\end{table}

The first three rows (italic) are lint counts computed directly from the
project's built-in v3 source by zlint's Source metadata field (unit:
lints; the 14 CRL lints are identified via RegisterRevocationListLint
and fall outside our single-cert scope); the remaining rows are the coverage
verdicts of our lintable rules (unit: rules). ``full'' counts ``whether one of
our rules is fully implemented by some cognate zlint lint''; it does not form a
simple ratio with the total zlint lint count, mainly because one lint can hit
multiple rules.

\subsection{Lintability External Validation: Rule Extraction and Lintability
Determination}
\label{subsec:lintability-ext}

The first half of the framework---rule extraction and lintability
determination---needs an independent external reference. For this we bring in
the CABF BR mapping table published by zlint
maintainers~\cite{zlint-mapping-sheet}: this table is manually annotated rule by rule
by zlint experts as to ``whether a BR requirement can be expressed by a static
lint,'' unaware of our extractor, IR, and lintability decider. Each Yes/No label
in the table corresponds exactly to our joint output ``extract this rule $+$
decide whether it is lintable'' for the same rule; agreement indicates that our
rule extraction and lintability determination match human experts. The table
covers two versions, BR~1.4.8~\cite{cabf-br148-md} and BR~2.0.2~\cite{cabf-br202-md}, compared version-for-version against
our CABF-Server-1.4.8 backend (422 rules) and CABF-Server-2.0.2
backend (1024 rules) extracted by the same two-layer
pipeline.

\begin{table}[tb]
\centering
\scriptsize
\setlength{\tabcolsep}{3pt}
\caption{Version-for-version external validation on BR~1.4.8 / BR~2.0.2.}
\label{tab:extval}
\begin{tabular}{@{}lrrrrr@{}}
\toprule
Version & Sheet rows (Yes/No) & Match cov. & Agreement &
  Cohen's $\kappa$ & F1 \\
\midrule
BR~1.4.8 & 122 (75/47) & 86.9\% & 91.5\% & 0.823 & 0.930 \\
BR~2.0.2 & 250 (17/233) & 86.8\% & 96.8\% & 0.616 & 0.632 \\
\bottomrule
\end{tabular}
\end{table}

A small number of sheet rows fall outside our scope---title-case or lowercase descriptors such as ``Required if'', ``Deprecated'', or lowercase ``shall'' that RFC~2119 processing correctly skips, and, on the BR~2.0.2 sheet, cross-references or chapter anchors with no matching backend clause---so match coverage is near 87\% on both versions. On BR~2.0.2 the matched positive class is very small (12 of 217 matched pairs, 5.5\%), so a naive all-No baseline already attains 94.5\% agreement; thus $\kappa = 0.616$ is the more meaningful figure, and the positive-class metrics ($P = 0.857$, $R = 0.50$, $n = 12$) carry correspondingly wide confidence intervals. BR~1.4.8 has more balanced classes (69 of 106 matched pairs positive, 65\%) and yields $\kappa = 0.823$ with positive-class $F_1 = 0.930$ and zero false positives, a tighter estimate of system quality.

This external validation checks the first half of the system. On
BR~1.4.8, $\kappa = 0.823$ and Yes recall 87.0\% indicate that our rule
extraction (identifying which of the 422/1024 backend items are genuine
normative requirements) and lintability determination (the four-condition
conjunction, \S\ref{sec:lintability}) agree strongly with zlint experts'
independent manual judgment; the 50.0\% Yes recall on BR~2.0.2 indicates a clear
false-negative margin remaining in the newer BR version. This validation does
not check downstream code generation and synonymy; those are checked
additionally on real certificates by the certificate detection of
\S\ref{subsec:certdetect}.

\subsection{Certificate Detection}
\label{subsec:certdetect}

Injecting the 90 final-shipping-strict lints into the zlint binary and scanning
certificates checks code generation and synonymy judgment. We count a hit as a
valid finding only when three conditions hold: the final zlint code's semantics
have been checked back to the source text, the certificate yields a non-pass
result, and an independent structural audit re-derives the same defect from the
DER / \texttt{openssl}; on external corpora we additionally require that no native
zlint lint flags the same certificate, so that existing coverage is not
miscredited to a generated lint. The unit is a (certificate $\times$ lint)
finding. The corpus first uses zlint's own testdata as a gate rather than an
ecosystem estimate: on 1128 parsable certificates the 90 lints raise 2627 raw
cicasgen\_ hits, of which 299 are independently confirmed (251 Error $+$
48 Warn) and 2328 are \texttt{NOCHECK} and hence not claimed.

\begin{table*}[tb]
\centering
\scriptsize
\setlength{\tabcolsep}{4pt}
\caption{Strictly confirmed findings across the three certificate corpora
(highest severity first; 90 shipped lints). ``ext.\ merged'' is the Tranco/CT
union after cross-corpus de-duplication.}
\label{tab:certfind}
\begin{tabular}{@{}l p{5.9cm} r r r r p{3.3cm}@{}}
\toprule
Sev. & Finding type / governing obligation & testdata & Tranco & CT &
  ext.\ merged & Evidence boundary \\
\midrule
Error & Subscriber-cert Certificate Policies must contain exactly one CABF
  Reserved Certificate Policy Identifier (BR~\S7.1.2.7.9) & 113 & 1 & 0 & 1 &
  one Tranco leaf lacks the reserved policy OID; no native zlint overlap \\
Error & AKI must not contain authorityCertSerialNumber / authorityCertIssuer
  (BR~\S7.1.2.11.1) & 32 & 2 & 0 & 2 & one Tranco CA cert violates both AKI
  subfields; no native zlint overlap \\
Error & Subscriber-cert Certificate Policies must not contain anyPolicy
  (BR~\S7.1.2.7.9) & 26 & 0 & 0 & 0 & confirmed only in testdata fixtures \\
Error & IV/OV subject-profile constraints on surname, givenName, and
  localityName / stateOrProvinceName (BR~\S7.1.2.7.3--\S7.1.2.7.4) & 57 & 0 & 0 &
  0 & testdata fixtures only \\
Error & Other structural constraints: SAN critical, signatureAlgorithm matches
  tbsCertificate.signature, subscriber pathLenConstraint forbidden, empty
  Certificate Policies, v1 uniqueID forbidden (RFC~5280 / BR) & 23 & 0 & 0 & 0 &
  testdata fixtures only \\
Warn & A certificate carrying only basic fields SHOULD be v1
  (RFC~5280~\S4.1.2.1) & 47 & 0 & 0 & 0 & advisory, reported as lint.Warn \\
Warn & Root CA Certificate SHOULD NOT contain CRL Distribution Points
  (BR~\S7.1.2.11.2) & 1 & 3 & 1 & 3 & advisory profile warning; the CT finding
  duplicates the same Root CA as one Tranco finding \\
\bottomrule
\end{tabular}
\end{table*}

The testdata column sums to 299 confirmed fixture findings (251 Error $+$ 48
Warn); Tranco contributes 6 and CT 1, i.e., 7 findings by source. We read
severity conservatively: the Certificate-Policies and AKI violations are
Error-level profile faults whose direct cryptographic consequence is indirect,
while the Root CA CRLDP finding is a lower, Warn-level advisory.

\textbf{Supplementary scan on real certificates.} To test whether these findings
exist only in hand-constructed testdata, we scan two external corpora: the
TLS-handshake certificate chains of the Tranco Top~1M domains (47,791
deduplicated PEMs) and a recent sample of CT logs (63,327 deduplicated PEMs). The
raw cicasgen\_ hit counts are large---44,020 on Tranco and 57,558 on CT---but
cannot be taken as defects: if a generated lint loses profile scope, a
conditional clause, or an effective-date constraint, an independent audit may
still only prove that the certificate satisfies the trigger of the wrong lint. We
therefore treat the external corpora only as negative validation. Under the
strict no-upstream audit these raw hits reduce to 6 confirmed on Tranco and 1 on
CT; summed by source that is 7 findings, and after cross-corpus de-duplication it
is 6 (certificate $\times$ lint) findings over 5 unique certificates, because one
Root CA CRLDP warning is confirmed on the same Comodo ``AAA Certificate
Services'' root in both corpora.

Raw hit rates therefore serve only to surface audit candidates, never as
defect-count conclusions. This supports the negative message of
Table~\ref{tab:snapshot}: even when a lint compiles, executes, and fires on a
certificate, one must still re-check the source-text scope, its conditional
clauses, and rule effective dates before calling it a defect. Concretely,
several high-frequency cicasgen\_ hits are excluded on exactly these
grounds---an OV/IV localityName requirement that dropped its
``stateOrProvinceName absent'' precondition (all such certs carry
stateOrProvinceName); an organizationalUnitName or AKI-subfield prohibition
applied to certificates issued before the requirement's effective date; the
RFC~5280 A.1 ``version MUST be v2'' fragment mis-extracted as a profile rule that
flags normal v3 certificates; and SubCA anyPolicy / NameConstraints rules that
approximate profile scope with IsSubCA alone. Hence the raw hit rate must not be
read as the real compliance-defect rate.

%-------------------------------------------------------------------------------
\section{Discussion}
\label{sec:discussion}
%-------------------------------------------------------------------------------

\subsection{Applicability Across PKIs}
\label{subsec:applicability}

The framework's core mechanism is not specific to PKI: Layer~1's RFC~2119
keyword set can be reconfigured to another ecosystem's deontic keywords (e.g.,
ISO's shall/should/may). The method applies to specification sources that
jointly satisfy three conditions: \textbf{(1)} they express mandatory/recommended
levels with recognizable formal deontic keywords; \textbf{(2)} the constrained
target is a structured, machine-parsable artifact (certificate, protocol
message, configuration file); and \textbf{(3)} they have a hierarchical section
structure that supports a knowledge graph and scope inheritance. This gives the
methodology value beyond PKI alone.

\subsection{Recommendations for Specification Authors and Tool Maintainers}
\label{subsec:recommendations}

For \textbf{specification authors}, three drafting principles reduce extraction
ambiguity: one obligation per sentence, referring to fields precisely by ASN.1
path rather than pronoun, and normalizing cross-document references with a fixed
pattern. For \textbf{lint maintainers}, we suggest writing the description
directly as a code\_summary that explicitly gives the field
path/assertion type/value/severity and cites the specification clause---we find
that many lint descriptions are neither specification citations nor faithful
implementation summaries, forcing the coverage analysis to first compute a
code\_summary before aligning, adding overhead and adjudication noise.

\subsection{Cross-Standard Coverage and Source Partitioning}
\label{subsec:crossstd}

The coverage analysis also exposes a scoping issue: of the 90 lints in the
emitted set, \textbf{9} have specification text from CABF BR yet are already
implemented by existing \textbf{RFC~5280} lints in zlint. A rule-by-rule check
against the zlint v3 source confirms the following.

\begin{table*}[tb]
\centering
\footnotesize
\caption{Nine emitted lints whose CABF BR text is already implemented by native
RFC~5280 lints in zlint.}
\label{tab:crossstd}
\begin{tabular}{@{}p{5.5cm}r p{9cm}@{}}
\toprule
CABF BR rule (source of generated lint) & \# & zlint lint that implements it \\
\midrule
issuerUniqueID / subjectUniqueID MUST NOT be present
  (7.1.2.1/2/8) & 6 &
  e\_cert\_contains\_unique\_identifier (RFC~5280:4.1.2.8) \\
signatureAlgorithm byte-for-byte matches
  tbsCertificate.signature (7.1.2.2) & 1 &
  e\_cert\_sig\_alg\_not\_match\_tbs\_sig\_alg (RFC~5280:4.1.1.2) \\
Subscriber/OCSP pathLenConstraint MUST NOT be present (7.1.2.7.8/8.4) &
  2 & e\_path\_len\_constraint\_improperly\_included
  (RFC~5280:4.2.1.9) \\
\bottomrule
\end{tabular}
\end{table*}

The cause is cross-standard inheritance of compliance scope: CABF BR~\S7.1.2
opens by declaring that its certificate profile ``incorporate, and are derived
from RFC~5280,'' and the normative requirements RFC~5280 imposes also apply. The
uniqueID prohibition, signatureAlgorithm consistency, and pathLenConstraint
restriction are originally RFC~5280 requirements, restated by CABF BR in
profile-table form. zlint implements the corresponding lints per RFC~5280 (their
Source field is RFC5280), but the coverage determination of
\S\ref{subsec:coverage} narrows candidates by rule source; a CABF BR rule is
compared only against lints with the CABF-BR source and never enters the
RFC~5280 lint candidate pool, so these 9 are judged none, wrongly enter the
codegen domain, and are eventually emitted.

\subsection{Limitations and Threats to Validity}
\label{subsec:limitations}

Our conclusions are strictly limited by the input representation used (structured
IR), the restricted code space, and the verification pipeline; they suit static
constraints that close within a single-document context and should not be
extrapolated to ``all PKI specifications can be fully automated.'' We note seven
boundaries.

\begin{itemize}
\item \textbf{(i) Method expressiveness.} The synonymy-judgment endpoint is still
  implemented by an LLM, and the 136 atomic templates still have gaps for
  byte-level encoding, host-unexposed fields, dynamic constraints, and the like.
\item \textbf{(ii) Component necessity not ablated.} The support for the
  knowledge graph / restricted DSL / SAIV is respectively an engineering
  trade-off, an architectural argument, and an un-ablated item; we do not claim
  quantitative superiority of graph retrieval over ordinary RAG.
\item \textbf{(iii) Missing modern-LLM baseline.} A direct-Go-generation
  comparison is not done, but certificate-level execution verification is
  independent of the specific generator, making it well-defined future work.
\item \textbf{(iv) Coverage and generalization boundary.} The lintable total 248
  and the codegen target 90 are different scopes, and for cross-ecosystem
  generalization we give only a mechanistic argument.
\item \textbf{(v) End-to-end scope.} ``Code $\equiv$ specification'' holds only
  for the emitted set (Table~\ref{tab:snapshot}, \S\ref{subsec:certdetect}), and
  the lintability label has not been exhaustively manually audited over all 248
  rules.
\item \textbf{(vi) Synonymy-endpoint blind spot.} ``code\_summary $\equiv$
  specification'' is given by an LLM judge comparing the code\_summary
  with the source text, not a model-independent check like certificate-level
  execution verification, and its ``source text'' is taken from the
  rule\_text recorded verbatim at extraction time (a field inside the
  IR) rather than re-read from the specification source; so if the original
  fragment or the IR has already dropped semantics, the synonymy judge may not
  independently detect it.
\item \textbf{(vii) Coverage determination partitions candidates by source.}
  \S\ref{subsec:coverage} retrieves only CABF-BR-source zlint lints for CABF BR
  rules, severing cross-standard inheritance of coverage: at least 9 rules
  originating in CABF BR but actually implemented by existing RFC~5280 lints in
  zlint are thereby misjudged as uncovered and wrongly enter the codegen domain
  (\S\ref{subsec:crossstd}). This underestimates the tools' true coverage and
  inflates the coverage-gap count. In addition, generated lints carry the
  cicasgen\_* prefix and are mutually rejected with zlint on both sides
  to prevent ``covering oneself.''
\end{itemize}

%-------------------------------------------------------------------------------
\section{Conclusion}
\label{sec:conclusion}
%-------------------------------------------------------------------------------

We presented an end-to-end framework that turns natural-language Web PKI
normative text into auditable static lint code. Applied to RFC~5280 and CABF~BR,
the pipeline recalled 2,080 candidate rules and identified 1,567 genuine
normative rules. Only 260 of these genuine rules, or 16.6\%, are lintable as
single-certificate static checks. Of these 260 lintable rules, existing zlint
lints fully cover only 165, leaving 95 uncovered and revealing a substantial gap
between deployed linting and the normative specifications.
For these uncovered rules, our restricted generator produced 77
compilable lints that contain at least one GENERIC atom and 16 further compilable
lints from registered NON\_GENERIC-only atom trees, for 93 compilable lints in
total. Of these, 79 passed the unanimous code--specification synonymy gate and
were injected into zlint. Across 111,118 deployed certificates, certificate
detection found four independently confirmed, cross-corpus de-duplicated
external certificate-problem instances that native zlint does not flag: one
Error-level AKI violation and three Warning-level Root CA CRLDP violations. We
hope our work can provide a useful reference for research related to certificate
issuance compliance automation.


\appendices
%-------------------------------------------------------------------------------
% Technical appendices A--F, English translation of the source appendices.
%-------------------------------------------------------------------------------

\section{Knowledge-Graph Relation Types and Retrieval Admissibility}
\label{app:kg}

The Deterministic Context Retrieval procedure of \S\ref{sec:method-graphrag} builds context around a set of
core relations and explicitly distinguishes ``relations admissible into
retrieval'' from ``relations used only in post-processing''
as shown in Table~\ref{tab:kgrel}. The retrieval stage
obeys four constraints, GR-1 through GR-4: no inferred normative rule may enter
retrieval; only raw specification nodes may enter retrieval; no inferential
relation may enter retrieval; and all context must be traceable to a
specification source.

\begin{table}[tb]
\centering
\footnotesize
\setlength{\tabcolsep}{4pt}
\caption{Knowledge-graph relation types and retrieval admissibility.}
\label{tab:kgrel}
\begin{tabularx}{\columnwidth}{@{}l l X c@{}}
\toprule
Type & Relation & Role & In retrieval \\
\midrule
core & CONTAINS & containment hierarchy among specs, sections, rules & Yes \\
core & DEFINES & link from a term/field/definition passage to its defined
  object & Yes \\
core & REFERENCES & explicit reference to another spec/section/algorithm & Yes \\
core & APPLIES TO & link from a rule to the abstract concept it constrains
  & Yes \\
aux & AFFECTS & link from a rule to the concrete certificate field it affects;
  used by the field resolver & No \\
aux & DERIVED FROM & provenance link from a derived rule to its source section
  & No \\
inferred & OVERRIDES & a more specific rule overrides a general one (produced by
  the rule engine) & No \\
inferred & CONFLICTS WITH & potential conflict between rules (produced by the
  rule engine) & No \\
\bottomrule
\end{tabularx}
\end{table}

\section{Key Fields of the Structured IR}
\label{app:ir}

The system internally uses a JSON-serialized, machine-readable IR object.
Beyond the core four-tuple of \S\ref{sec:method-dual-layer}, its key fields are listed in Table~\ref{tab:irfields};
the remaining fields support provenance, normalization, matching, and downstream
product generation. The lintability decision in \S\ref{sec:lintability} directly
uses only the four fields obligation, assertion\_subject, enforcement\_phase, and
rule\_category.

\begin{table}[tb]
\centering
\footnotesize
\caption{Key fields of the structured IR.}
\label{tab:irfields}
\begin{tabularx}{\columnwidth}{@{}l l X@{}}
\toprule
Field & Role & Meaning \\
\midrule
subject & core & constrained certificate field or path \\
obligation & core & RFC~2119 deontic level, such as MUST/SHOULD/\dots \\
predicate & core & constraint relation, such as presence, equality,
  containment, or range \\
constraint & core & constraint value, pattern, or condition \\
rule\_category & key ext. & rule's semantic category, such as
  encoding\_constraint, definition, algorithm\_ref, \dots \\
assertion\_subject & key ext. & assertion subject, such as Certificate / CA /
  RelyingParty / external ecosystem \\
enforcement\_phase & key ext. & phase the constraint depends on
  such as encoding / runtime / external validation \\
source\_section & key ext. & specification source and section
  identifier \\
source\_span & key ext. & character/sentence span in the source text \\
evidence\_text & key ext. & source fragment supporting this IR record \\
context\_nodes & key ext. & definition, field metadata, or cited
  context nodes from Deterministic Context Retrieval \\
\bottomrule
\end{tabularx}
\end{table}

\section{Construction Details of the Restricted Code Space $\Tv$}
\label{app:tv}

\textbf{The seven components of the vocabulary $\mathcal{V}$.} $\mathcal{V}$ is
the disjoint union of seven finite sets
shown in Table~\ref{tab:vocab}: certificate fields $\mathcal{F}_{\mathrm{cert}}$,
DN fields $\mathcal{F}_{\mathrm{dn}}$ over the Subject and Issuer names, OID
constants $\mathcal{O}$, including zlint util.* extension and algorithm OIDs
together with well-known CABF and RFC OIDs, KeyUsage bits $\mathcal{B}_{\mathrm{KU}}$,
ExtKeyUsage bits $\mathcal{B}_{\mathrm{EKU}}$, ASN.1 encoding types
$\mathcal{E}_{\mathrm{ASN1}}$, and named regexes $\mathcal{R}_{\mathrm{regex}}$,
each an audited literal RE2 pattern. Table~\ref{tab:vocab} gives the size and
representative members of each component.
All components are finite sets; the LLM receives a full enumeration of
each component in the prompt, so any identifier outside $\mathcal{V}$ in its
output triggers an $\eta$ parse error.

\begin{table*}[tb]
\centering
\footnotesize
\caption{The seven frozen components of the restricted vocabulary $\mathcal{V}$,
with their sizes and representative members. The LLM may emit only identifiers
drawn from these closed sets; the full enumerations are released with the code
and data.}
\label{tab:vocab}
\begin{tabularx}{\textwidth}{@{}l l r X@{}}
\toprule
Symbol & Component & Size & Representative members \\
\midrule
$\mathcal{F}_{\mathrm{cert}}$ & Certificate fields & 52 &
  c.Version, c.NotAfter, c.KeyUsage, c.DNSNames, c.IPAddresses \\
$\mathcal{F}_{\mathrm{dn}}$ & DN fields over Subject / Issuer & 38 &
  Subject.CommonName, Subject.Organization, Subject.Country, Issuer.CommonName \\
$\mathcal{O}$ & OID constants & 99 &
  util.BasicConstOID, util.SubjectAlternateNameOID, util.AiaOID \\
$\mathcal{B}_{\mathrm{KU}}$ & KeyUsage bits & 9 &
  DigitalSignature, KeyEncipherment, CertSign, CRLSign \\
$\mathcal{B}_{\mathrm{EKU}}$ & ExtKeyUsage bits & 11 &
  ServerAuth, ClientAuth, CodeSigning, EmailProtection \\
$\mathcal{E}_{\mathrm{ASN1}}$ & ASN.1 encoding types & 11 &
  PrintableString, IA5String, UTF8String, BMPString \\
$\mathcal{R}_{\mathrm{regex}}$ & Named RE2 regexes & 28 &
  Re\_LDH\_Hostname, Re\_Rfc5321Mailbox, Re\_TorV3Onion \\
\bottomrule
\end{tabularx}
\end{table*}

\textbf{Representative atomic templates.} The atomic-template set $\mathcal{A}$
has $|\mathcal{A}|=151$ and is organized by semantic cluster, each with a typed
signature $\mathrm{sig}(a)$ described in \S\ref{subsec:tv}. Table~\ref{tab:atoms} excerpts
key atomic templates by cluster, and all 151 are released with the code and data.

\begin{table*}[tb]
\centering
\footnotesize
\caption{Representative atomic templates by semantic cluster.}
\label{tab:atoms}
\begin{tabularx}{\textwidth}{@{}l p{4.2cm} l X@{}}
\toprule
Cluster & Atomic template & Signature & Semantics \\
\midrule
I Extension presence & ExtPresent / ExtCritical &
  $(\mathcal{O})$ & extension OID present / present and critical bit set \\
II Macro attributes & IsCA / KeyUsageHas /
  ExtKeyUsageHas & $()$ / $(\mathcal{B}_{\mathrm{KU}})$ /
  $(\mathcal{B}_{\mathrm{EKU}})$ & BasicConstraints.cA$=$true / KeyUsage bit /
  EKU item \\
III Field value/form & FieldEq / FieldMatchesRegex /
  FieldEncodedAs & $(\mathcal{F}, \cdot)$ & field equals literal /
  matches named regex / ASN.1 encoding type \\
IV List/byte-level & ListAllMatch / OidListContains /
  BytesContainsOidDer & $(\mathcal{F}, \cdot)$ & list all satisfy
  subtree / contains named OID / bytes contain the OID's DER \\
V NameConstraints & SubtreeIPListAny\-HasOctetCount\-AndNotAllZero &
  $(\mathcal{F}, \mathbb{Z})$ & NC IP subtree has an entry of given octet count
  and not all-zero \\
VI Special structure & DomainComponentOrdered & $()$ &
  domainComponent in reverse order per RFC~4519 \\
\bottomrule
\end{tabularx}
\end{table*}

There are three combinators $\{\neg, \wedge, \vee\}$, with semantics following
classical propositional logic.

\textbf{Generality grading of atomic templates: GENERIC and NON\_GENERIC.} Atomic
templates are graded by two orthogonal criteria independently of the number of
parameters: an atomic template is GENERIC if and only if, first, it expresses a
general PKI concept---a class of certificate-attribute judgments reusable
across rules---and, second, it is parameterized over fields/values and bound to no
specific rule\_id or corpus-specific OID/single clause; otherwise it is
NON\_GENERIC. In that case, its logic is specialized to the internal structure of
some extension or a single RFC/CABF clause, and its semantics remain fixed to
that construct even when parameterized. A zero-parameter atomic template may
belong to either class. For example, IsCA is GENERIC, while
NotAfterIsNoExpirySentinel is NON\_GENERIC.
This grading is fixed in code as the two sets GENERIC\_ATOMS /
NON\_GENERIC\_ATOMS with a partition-completeness assertion, and is
deterministically recomputable.

Of the 151 atomic templates in $\mathcal{A}$, 98 are GENERIC and 53 are
NON\_GENERIC. Table~\ref{tab:nongeneric} lists only representative NON\_GENERIC
examples; the full list is released with the code and data.

\begin{table*}[tb]
\centering
\footnotesize
\caption{Representative NON\_GENERIC atomic templates.}
\label{tab:nongeneric}
\begin{tabularx}{\textwidth}{@{}p{4.6cm} p{3cm} X@{}}
\toprule
Representative NON\_GENERIC atomic template & Specialization target &
  Semantics \\
\midrule
SigAlgMatchesTBSSignature & certificate signature-algorithm fields &
  outer signatureAlgorithm byte-for-byte identical to
  tbsCertificate.signature \\
NotAfterIsNoExpirySentinel & validity sentinel value & whether
  notAfter is a specific no-expiry sentinel time \\
DomainComponentOrdered & Subject DN's domainComponent &
  whether the domainComponent sequence is ordered per RFC~4519 \\
CertPolicy\-ExplicitText\-HasEncodingTagInSet & certificatePolicies
  explicitText & whether explicitText uses an allowed ASN.1 string encoding \\
CRLDPHasNameRelative & CRL Distribution Points & whether
  distributionPoint uses the nameRelativeToCRLIssuer form \\
SubtreeIPListAny\-HasOctetCount\-AndNotAllZero & NameConstraints IP subtree
  & whether the IP subtree has an entry of given octet count that is not
  all-zero \\
WildcardFilter & DNS name patterns & whether a domain wildcard
  satisfies the specialized filter rule \\
\bottomrule
\end{tabularx}
\end{table*}

NON\_GENERIC atomic templates are allowed into code generation---they correspond
to real PKI constructs, not padding to inflate the corpus---but are explicitly
marked. The GENERIC/NON\_GENERIC split of the 93 generated lints is reported in
\S\ref{sec:eval}; among the 79 final-shipping unanimous lints the split is 60
GENERIC-only, 4 mixed, and 15 NON\_GENERIC-only.

\textbf{Rendering and metadata binding.} After $\rho$ renders them, all DSL trees
are embedded into a fixed zlint Go shell, with registration via
RegisterCertificateLint and the CheckApplies/Execute
method pair, so inter-rule differences are entirely concentrated in the check
body and imports; $\Phi_{\mathrm{post}}$ of \S\ref{subsec:synth} deterministically
binds Description/Citation/Name and each source's
PACKAGE/SOURCE/EFFECTIVE\_DATE metadata. For example,
RFC5280 $\mapsto$ RFC5280, CABF-TLS-BR $\mapsto$
CABFBaselineRequirements. The obligation-to-severity mapping is
MUST/MUST NOT/SHALL/SHALL NOT/REQUIRED $\mapsto$ lint.Error with lint-name
prefix e\_, and SHOULD/SHOULD NOT/RECOMMENDED $\mapsto$
lint.Warn with lint-name prefix w\_; MAY/OPTIONAL is already excluded by $C_1$ of
\S\ref{sec:lintability}, so the system produces no lint.Notice-level
output.

\section{Phrase Dictionary of the Mechanical Translation $\smech$}
\label{app:sigma}

$\smech$, defined in \S\ref{subsec:mech}, consists of an
atomic-template-to-phrase dictionary $\mathcal{M} : \mathcal{A} \to
\mathcal{L}_{\mathrm{NL}}$ and three combinator reduction rules. Representative
entries appear in Table~\ref{tab:phrasedict}.

\begin{table}[tb]
\centering
\footnotesize
\caption{Representative entries of the phrase dictionary $\mathcal{M}$.}
\label{tab:phrasedict}
\begin{tabularx}{\columnwidth}{@{}p{3.1cm} X@{}}
\toprule
Atomic template & $\mathcal{M}(a)$ phrase template \\
\midrule
ExtPresent(O) & ``the \{O\} extension is present'' \\
ExtCritical(O) & ``the \{O\} extension is present and marked
  critical'' \\
FieldEq(F, v) & ``\{F\} equals \{v\}'' \\
FieldEncodedAs(F, T) & ``\{F\} is encoded as ASN.1 \{T\}'' \\
KeyUsageHas(B) & ``the KeyUsage bit \{B\} is asserted'' \\
ListAllMatch(F, T) & ``every entry of \{F\} satisfies
  the result of $\smech(T)$'' \\
SubtreeIPListAny\-HasOctetCount\-AndNotAllZero(F, n) & ``the
  NameConstraints IP subtree \{F\} contains a non-zero entry of \{n\}
  octets'' \\
\bottomrule
\end{tabularx}
\end{table}

Combinator reduction: $\smech(\neg t) = $ ``NOT ($\smech(t)$)'';
$\smech(t_1 \wedge t_2) = $ ``($\smech(t_1)$) AND ($\smech(t_2)$)''; $\vee$
likewise; the atomic-template base case $\smech(a) = \mathcal{M}(a)$. For a
conditional pair $(p,q)$: when $p = \neg(p')$, it outputs ``WHEN NOT
($\smech(p')$), THEN $\smech(q)$'', the NEG-PRE template that eliminates double
negation; otherwise it outputs ``WHEN ($\smech(p)$), THEN $\smech(q)$''; and
when $p = \varepsilon$ it outputs the single sentence $\smech(q)$.

%-------------------------------------------------------------------------------
% Stage-Aware Iterative Verification, abbreviated SAIV --- moved out of the main body and
% merged inline here; previously the separate file 07_saiv.tex, now deleted.
%-------------------------------------------------------------------------------
\section{Stage-Aware Iterative Verification: SAIV}
\label{sec:saiv}

To coordinate the modules and safeguard their output quality without manual
ground truth, we propose the Stage-Aware Iterative Verification framework,
abbreviated SAIV.
The basic idea is to predefine two quantifiable optimization
objectives, measure the gap between the current and ideal states by each
objective's computable residual, and when a residual has not reached its
threshold, the operator uses the residual signal to locate the offending module
and applies the corresponding repair, then re-measures and proceeds to the next
iteration, until both residuals converge. Alongside the two objectives, two
auxiliary checks police the lintability decision without entering this loop.

\subsection{Formal Definition}
\label{subsec:saiv-formal}

The whole pipeline $\Pi$ chains four modules---extraction $\to$ classification
$\to$ generation $\to$ verification:
\begin{equation}\label{eq:pipeline}
  \Pi = \phi_V \circ \phi_G \circ \phi_C \circ \phi_R
\end{equation}
where $\phi_R$ is extraction, $\phi_C$ is the four-condition lintability decision
of \S\ref{sec:lintability}, $\phi_G$ is the DSL code space $\Tv$ of
\S\ref{sec:codegen}, and $\phi_V$ is the three-layer semantic alignment of
\S\ref{sec:codegen}. $\phi_C$ is the
single-rule lintability predicate $r \mapsto \{0,1\}$; the $\psi_C$ of
Algorithm~\ref{alg:saiv} is its set-level lift, which strips the noise $\RN$ from
$\Rkw$ and sorts the remaining genuine rules into lintable $\RL$ and not-lintable
$\RU$, giving the partition $(\RN, \RL, \RU)$. The operator drives the loop from
the computable signals these modules produce.

\subsection{Two Iteration Objectives and Two Auxiliary Lintability Checks}
\label{subsec:targets}

SAIV borrows the idea of backpropagation. It fixes two optimization objectives
the whole pipeline is trained to approach, and keeps apart from them two
auxiliary checks that police the lintability decision $\phi_C$ rather than drive
optimization. The two objectives G1 and G2 each carry a directly computable
residual, $\Lrec$ and $\Lcode$, whose weighted sum is the total loss $\Ltot$, so
convergence is defined by $\Ltot$ alone. S1 and S2 are checks, not objectives,
and add no term to $\Ltot$. They form the two sides of a single consistency test
on $\phi_C$, S1 auditing the not-lintable verdict and S2 the lintable verdict,
each firing only to correct a suspected misclassification. The subsequent subsections
revolve around the two objectives, while the two checks gate the labels those
objectives consume.

\medskip
\noindent\textbf{G1: recall completeness.}
\begin{theorem}[Recall-Completeness Invariant]\label{thm:recall}
If the classification $\psi_C$ assigns each keyword-recalled candidate in
$\Rkw(\mathcal{D})$ to exactly one of the three classes---noise, lintable, and
not-lintable---then recall is conserved from recall to classification:
\begin{equation}\label{eq:cons}
  |\Rkw(\mathcal{D})| = |\RN| + |\RL| + |\RU|
\end{equation}
\end{theorem}
\noindent where $\mathcal{D}$ is the standard-document corpus, written D in
ASCII inside the Algorithm~\ref{alg:saiv} listing, $\RN$ is noise that contains
RFC~2119 keywords but not genuine rules, $\RL$ the lintable rules, and $\RU$
the not-lintable rules. The executable class $\RL$ is further split by ``whether
it is already implemented by an existing external lint tool'':
\begin{equation}\label{eq:split}
  |\RL| = |\RLcov| + |\RLunc|
\end{equation}
where $\RLcov$ are rules covered by some lint tool and $\RLunc$ are
not-yet-covered rules. The first identity~\eqref{eq:cons} is total conservation
from recall to classification, and the second~\eqref{eq:split} splits the
executable set so the later objectives can be pursued separately. The uncovered
subset $\RLunc$ is the most important observation domain for generation that
``fills the ecosystem gap,'' while the covered subset $\RLcov$ provides external
evidence of reverse executability. The residual is the recall-conservation
residual $\Lrec$, a symmetric normalization of the gap in
Theorem~\ref{thm:recall}, computable without manual ground truth and currently
closed, as shown in Table~\ref{tab:snapshot}:
\begin{equation}\label{eq:recall-res}
  \Lrec = 1 - \frac{\min\!\big(|\Rkw|,\; |\RN|+|\RL|+|\RU|\big)}
  {\max\!\big(|\Rkw|,\; |\RN|+|\RL|+|\RU|\big)}
\end{equation}

\medskip
\noindent\textbf{G2: code--specification synonymy for uncovered rules.} For the
uncovered subset $\RLunc$, the code summary $\smech(t_r)$, the code\_summary of
rule $r$'s DSL tree $t_r$ computed by $\smech$ of \S\ref{subsec:mech}, is
required to be synonymous with the specification text.
The residual $\Lcode$ is the average, over $\RLunc$, of the per-rule correctness
label $\lambda_{\mathrm{code}}(r) = \ind[\mathrm{compile}(\rho(t_r))] \cdot
\mathrm{syn}(\smech(t_r), \mathrm{spec}(r))$, where $\rho(t_r)$ is the rendered Go
code and $\mathrm{spec}(r)$ the rule text. The structural-completeness score is
$\equiv 1$ in the fixed zlint shell, so it drops out:
\begin{equation}\label{eq:code-res}
  \Lcode = 1 - \frac{1}{|\RLunc|}\sum_{r\in\RLunc}\lambda_{\mathrm{code}}(r)
\end{equation}

\medskip
\noindent\textbf{S1: reverse executability of covered rules.} The first
auxiliary check, auditing the not-lintable verdict. If a rule has a corresponding
implementation in a lint tool such as zlint it must be lintable, so the check
counts the violations $\Nviol$ and is clear when $\Nviol = 0$. The operator
computes $\Nviol$ by a reverse coverage check---running the coverage judge over
the not-lintable set $\RU$ and counting the rules it rates full or partial; each
such hit is adjudicated as either a $\phi_C$ false negative and re-extracted to
lintable, or a spurious coverage match and kept not-lintable.

\medskip
\noindent\textbf{S2: synthesizability of lintable rules.} The dual check,
auditing the lintable verdict. A rule judged lintable should be synthesizable
within $\Tv$, so a no-template abstention $\bot_{\mathrm{NT}}$ returned for such a
rule after generation repair $\mathcal{P}_G$ and vocabulary expansion
$\mathcal{P}_A$ are exhausted is a candidate signal that $\phi_C$ mislabeled it.
S2 adds no residual to $\Ltot$; it reuses the generator's existing abstention
marker as a diagnostic, and each persistent abstention is reviewed against the
four-condition decision of \S\ref{sec:lintability} and adjudicated as either a
$\phi_C$ false positive and re-derived to not-lintable, or a genuine $\Tv$
expressiveness gap and kept as a logged abstention. In the shipping corpus both
no-template abstentions are $\Tv$ gaps, so neither is re-adjudicated as
not-lintable.

\medskip
The two objective residuals are weighted into the total loss in
Eq.~\eqref{eq:total}, with $w_R + w_C = 1$:
\begin{equation}\label{eq:total}
  \Ltot = w_R\cdot\Lrec + w_C\cdot\Lcode
\end{equation}
Since recall is conserved, $\Lrec = 0$, $\Ltot$ is governed by $\Lcode$ alone.
The two auxiliary checks stay outside $\Ltot$; their signals are the violation
count $\Nviol$ of S1 and the abstention review of S2 in \S\ref{subsec:algo}.

\begin{table*}[tb]
\centering
\footnotesize
\caption{The two SAIV objectives and the two auxiliary lintability checks, with
their signals and repair routing.}
\label{tab:saiv-objectives}
\begin{tabularx}{\textwidth}{@{}llp{4.0cm}X@{}}
\toprule
Item & Signal & Main repair direction & Repair target \\
\midrule
\multicolumn{4}{@{}l}{\emph{Iteration objectives --- define $\Ltot$ and
  convergence}}\\
G1 Recall completeness & $\Lrec$, Eq.~\eqref{eq:recall-res} & --- &
  recall logic \\
G2 Synonymy & $\Lcode$, Eq.~\eqref{eq:code-res} &
  IR-content, DSL-generation, and verification-judgment repairs &
  IR-content correction / DSL-tree re-synthesis / judgment \& summary
  re-check \\
\midrule
\multicolumn{4}{@{}l}{\emph{Auxiliary lintability checks --- police $\phi_C$,
  outside $\Ltot$}}\\
S1 Reverse executability & $\Nviol$, \S\ref{subsec:algo} &
  verification-judgment consistency repair &
  $\phi_C$ false negatives / spurious coverage matches \\
S2 Synthesizability & $\bot_{\mathrm{NT}}$ abstention, \S\ref{subsec:algo} &
  generation repair then vocabulary expansion &
  $\phi_C$ false positives / $\Tv$ expressiveness gaps \\
\bottomrule
\end{tabularx}
\end{table*}

Table~\ref{tab:saiv-objectives} gives only the item-to-repair-direction correspondence; the
concrete repair operations and their notation are defined in
\S\ref{subsec:repair}.

\subsection{Stage Routing and Repair Operations}
\label{subsec:repair}

\textbf{Stage routing.} Each round the operator chooses the next repair operation using
directly computable quantities, namely the recall-conservation residual, the
compilation-failure rate over the lintable set, the average structural score,
the entity-necessary-condition screening result, and the synonymy confidence.
This routing only decides which kind of repair to attempt next round.

\textbf{Repair operations.} SAIV comprises four kinds of operations---IR-content
repair $\mathcal{P}_R$, classification repair $\mathcal{P}_C$, generation repair
$\mathcal{P}_G$, and verification repair $\mathcal{P}_V$, the subscripts naming
the repaired module $\phi_\bullet$. The repair operator $\mathcal{P}$ is distinct
from the rendering function $\rho$ of \S\ref{subsec:mech}. The operator selects
the next round's operation from the signals above, and a repair is accepted only
when the metric decreases.

\begin{itemize}
\item \textbf{IR-content repair $\mathcal{P}_R$}: feeds the failure trace---the
  current IR, the irreducible category, the mechanical summary, the judge ruling
  and rationale, the synonymy confidence, and the session's IR history---back to
  the LLM, which returns REPAIR(ir') or NO\_FIX.
\item \textbf{Classification repair $\mathcal{P}_C$}: re-extracts the
  four-condition IR fields and re-runs the lintability decision.
\item \textbf{Generation repair $\mathcal{P}_G$}: tries deterministic repair
  first; on failure it feeds the failing code fragment and the IR-constraint
  discrepancy, or the parser's signature/closure error, back into the prompt to
  re-synthesize within $\Tv$; on persistent abstention it enters the offline
  vocabulary-expansion channel $\mathcal{P}_A$, which hand-authors a new
  atomic-template candidate and merges it into $\mathcal{A}$ after
  certificate-level execution verification.
\item \textbf{Verification repair $\mathcal{P}_V$}: enlarges the synonymy
  judgment's semantic neighborhood, e.g., negation/affirmation interchange or
  one-to-many decomposition, or corrects the inconsistency between the two
  verdict sources.
\end{itemize}

\subsection{Iteration Algorithm and Termination Conditions}
\label{subsec:algo}

Algorithm~\ref{alg:saiv} gives the full procedure.

% Algorithm 2 -- Stage-Aware Iterative Verification, abbreviated SAIV.
% Included from the SAIV section of appendix.tex via % Algorithm 2 -- Stage-Aware Iterative Verification, abbreviated SAIV.
% Included from the SAIV section of appendix.tex via \input{algorithms/alg_saiv}.
%
% Typeset single-column (not figure*) at \scriptsize, for the same reason as
% Algorithm 1: short algorithmic lines wasted the right half of a full-width float,
% which is billed two columns of page area. The RouteRepair \Comment is on its own
% \Statex line because \Comment is right-aligned and cannot wrap. No step, symbol,
% or comment wording was changed.
\begin{figure}[tb]
\setlength{\parskip}{0pt}
\makeatletter
\@ifundefined{c@algorithm}{%
  \newcounter{algorithm}%
  \renewcommand{\thealgorithm}{\arabic{algorithm}}%
}{}%
\makeatother
\refstepcounter{algorithm}
\label{alg:saiv}
\hrule width \linewidth height 0.6pt\relax
\vspace{1pt}
\nointerlineskip
\noindent\textbf{Algorithm~\thealgorithm.} Algorithmic listing for
Stage-Aware Iterative Verification, abbreviated SAIV.
\par
\vspace{1pt}
\nointerlineskip
\hrule width \linewidth height 0.4pt\relax
\nointerlineskip
\scriptsize
\begin{algorithmic}[1]
\Statex \textbf{Input:} Web PKI normative-source corpus $D$, threshold $\theta$, max iterations $K$
\Statex \textbf{Output:} final code set $C^{*}$ and termination status
\State $\Rkw \gets \phi_R(D)$;\quad $(\RN,\RL,\RU)\gets \psi_L(\Rkw)$
  \Comment{$\RU=$ not-lintable}
\State $(\RLcov,\RLunc)\gets \textsc{Coverage}(\RL)$
  \Comment{Algorithm~\ref{alg:coverage}}
\State $C \gets \{\phi_G(r): r\in \RLunc\}$
\State $t\gets 0$
\Repeat
  \State compute objective residuals $\Lrec$~\eqref{eq:recall-res},
    $\Lcode$~\eqref{eq:code-res}; lintability-check signals $\Nviol$ (S1) and the
    pending-abstention set $\NTpend$ (S2)
  \If{$\Ltot<\theta$ \textbf{and} $\Nviol=0$ \textbf{and} $|\NTpend|=0$}
    \State \textbf{break} \Comment{objectives converged, S1 and S2 clear}
  \EndIf
  \State $\mathit{op}\gets \textsc{RouteRepair}(\Lrec,\Lcode,\Nviol,\NTpend,\cdots)$
  \Statex \Comment{repair routing, \S\ref{subsec:repair}}
  \If{$\mathit{op}=\mathcal{P}_R$} \Comment{IR-content repair}
    \ForAll{$r$ routed to $\mathcal{P}_R$}
      \State $\mathit{rep}\gets \mathcal{P}_R(\textsc{FailTrace}(r))$
        \Comment{REPAIR(ir$'$) or NO\_FIX}
      \If{$\mathit{rep}=\textsc{No\_Fix}$ \textbf{or}
          \textsc{ReextractCheck}$(\mathit{rep}.ir',\textsc{text}(r))\neq\emptyset$}
        \State record $r$ in irreducible set $\mathcal{R}_{\mathrm{irred}}$
          \Comment{kept as residual}
      \ElsIf{\textsc{Retest}$(\mathit{rep}.ir')$ passes}
        \State $\mathrm{IR}(r)\gets \mathit{rep}.ir'$ \Comment{accept}
      \EndIf
    \EndFor
  \Else
    \State apply $\mathit{op}\in\{\mathcal{P}_L,\mathcal{P}_G,\mathcal{P}_V\}$
  \EndIf
  \State update $(\RN,\RL,\RU,\RLunc)$ and $C$
  \State $t\gets t+1$
\Until{$t\ge K$ \textbf{or} no residual decreased, the no-progress condition}
\State \textbf{return} $(C,\ \mathit{status})$, with $\mathit{status}=\textsc{closed}$
  on break, else \textsc{max\_iter} if $t\ge K$ and \textsc{dry} otherwise
\end{algorithmic}
\vspace{1pt}
\hrule width \linewidth height 0.6pt\relax
\end{figure}
.
%
% Typeset single-column (not figure*) at \scriptsize, for the same reason as
% Algorithm 1: short algorithmic lines wasted the right half of a full-width float,
% which is billed two columns of page area. The RouteRepair \Comment is on its own
% \Statex line because \Comment is right-aligned and cannot wrap. No step, symbol,
% or comment wording was changed.
\begin{figure}[tb]
\setlength{\parskip}{0pt}
\makeatletter
\@ifundefined{c@algorithm}{%
  \newcounter{algorithm}%
  \renewcommand{\thealgorithm}{\arabic{algorithm}}%
}{}%
\makeatother
\refstepcounter{algorithm}
\label{alg:saiv}
\hrule width \linewidth height 0.6pt\relax
\vspace{1pt}
\nointerlineskip
\noindent\textbf{Algorithm~\thealgorithm.} Algorithmic listing for
Stage-Aware Iterative Verification, abbreviated SAIV.
\par
\vspace{1pt}
\nointerlineskip
\hrule width \linewidth height 0.4pt\relax
\nointerlineskip
\scriptsize
\begin{algorithmic}[1]
\Statex \textbf{Input:} Web PKI normative-source corpus $D$, threshold $\theta$, max iterations $K$
\Statex \textbf{Output:} final code set $C^{*}$ and termination status
\State $\Rkw \gets \phi_R(D)$;\quad $(\RN,\RL,\RU)\gets \psi_L(\Rkw)$
  \Comment{$\RU=$ not-lintable}
\State $(\RLcov,\RLunc)\gets \textsc{Coverage}(\RL)$
  \Comment{Algorithm~\ref{alg:coverage}}
\State $C \gets \{\phi_G(r): r\in \RLunc\}$
\State $t\gets 0$
\Repeat
  \State compute objective residuals $\Lrec$~\eqref{eq:recall-res},
    $\Lcode$~\eqref{eq:code-res}; lintability-check signals $\Nviol$ (S1) and the
    pending-abstention set $\NTpend$ (S2)
  \If{$\Ltot<\theta$ \textbf{and} $\Nviol=0$ \textbf{and} $|\NTpend|=0$}
    \State \textbf{break} \Comment{objectives converged, S1 and S2 clear}
  \EndIf
  \State $\mathit{op}\gets \textsc{RouteRepair}(\Lrec,\Lcode,\Nviol,\NTpend,\cdots)$
  \Statex \Comment{repair routing, \S\ref{subsec:repair}}
  \If{$\mathit{op}=\mathcal{P}_R$} \Comment{IR-content repair}
    \ForAll{$r$ routed to $\mathcal{P}_R$}
      \State $\mathit{rep}\gets \mathcal{P}_R(\textsc{FailTrace}(r))$
        \Comment{REPAIR(ir$'$) or NO\_FIX}
      \If{$\mathit{rep}=\textsc{No\_Fix}$ \textbf{or}
          \textsc{ReextractCheck}$(\mathit{rep}.ir',\textsc{text}(r))\neq\emptyset$}
        \State record $r$ in irreducible set $\mathcal{R}_{\mathrm{irred}}$
          \Comment{kept as residual}
      \ElsIf{\textsc{Retest}$(\mathit{rep}.ir')$ passes}
        \State $\mathrm{IR}(r)\gets \mathit{rep}.ir'$ \Comment{accept}
      \EndIf
    \EndFor
  \Else
    \State apply $\mathit{op}\in\{\mathcal{P}_L,\mathcal{P}_G,\mathcal{P}_V\}$
  \EndIf
  \State update $(\RN,\RL,\RU,\RLunc)$ and $C$
  \State $t\gets t+1$
\Until{$t\ge K$ \textbf{or} no residual decreased, the no-progress condition}
\State \textbf{return} $(C,\ \mathit{status})$, with $\mathit{status}=\textsc{closed}$
  on break, else \textsc{max\_iter} if $t\ge K$ and \textsc{dry} otherwise
\end{algorithmic}
\vspace{1pt}
\hrule width \linewidth height 0.6pt\relax
\end{figure}


\textbf{Termination conditions}: $\Ltot < \theta$, with default $0.05$, and the two
lintability checks clear, reaching the maximum iteration $K = 10$, or no residual
decreasing in a round. If every
accepted stage-repair operation is monotone, reducing that stage's error or
leaving it unchanged, then $\Ltot$ is non-increasing each round; the actual
system uses a ``stop when no decrease'' policy, so it guarantees finite-round
termination but not global optimality, as discussed in \S\ref{sec:eval}.


\bibliographystyle{IEEEtran}
\bibliography{references}

%%%%%%%%%%%%%%%%%%%%%%%%%%%%%%%%%%%%%%%%%%%%%%%%%%%%%%%%%%%%%%%%%%%%%%%%%%%%%%%%
\end{document}
%%%%%%%%%%%%%%%%%%%%%%%%%%%%%%%%%%%%%%%%%%%%%%%%%%%%%%%%%%%%%%%%%%%%%%%%%%%%%%%%
